\subsection{\texorpdfstring{\(\mathbf{S}(\mathbf{M}^{*})\)、\(\mathbf{TS}(\mathbf{M})^{*\mathrm{gr}}\) 与 \(\operatorname{Pol}(\mathbf{M}, \mathbf{A})\) 之间的关系}{S(M*)、TS(M) 的分次对偶与 Pol(M, A) 之间的关系}}

设 M 为自由 A-模。我们将赋予分次对偶 \(\mathbf{TS}(\mathbf{M})^{*gr}\) 一个交换、结合且含单位元的分次 \(^{1}\) 代数结构，它由 \(\mathbf{TS}(\mathbf{M})\) 的分次余代数结构（III，第580页）导出。根据III，第497页，存在唯一的分次 A-代数同态

\[
\theta : \mathbf {S} (\mathbf {M} ^ {*}) \rightarrow \mathbf {T S} (\mathbf {M}) ^ {* \mathrm{gr}}
\]

它在次数 1 上诱导 \(\mathbf{M}^*\) 的恒等映射。

命题20. --- 若A-模M是自由且有限生成的，则 \(\theta\) 是分次代数的一个同构。

设 \((e_i)_{i\in I}\) 是 \(\mathbf{M}\) 的一组基，且 \((e_i^*)_{i\in I}\) 是 \(\mathbf{M}^*\) 的对偶基。对 \(\nu \in \mathbf{N}^I\) ，令

\[
e _ {\nu} = \prod_ {i \in I} \gamma_ {\nu_ {i}} (e _ {i}) \in \mathbf {T S} (\mathbf {M}).
\]

由命题4（IV，第47页），族 \((e_{\nu})_{\nu \in \mathbf{N}^{I}}\) 是 \(\mathbf{TS}(\mathbf{M})\) 的一组基；设 \((e_{\nu}^{*})\) 是 \(\mathbf{TS}(\mathbf{M})^{*\mathrm{gr}}\) 中与 \((e_{\nu})\) 对偶的基。根据III，第505页，定理1，只需证明对任意 \(\nu \in \mathbf{N}^{I}\) 有

\[
e _ {\nu} ^ {*} = \prod_ {i \in I} \left(e _ {i} ^ {*}\right) ^ {\nu_ {i}},
\]

或等价地，对 \(\rho, \sigma \in \mathbf{N}^{\mathrm{I}}\) 有 \(e_{\rho}^{*} \cdot e_{\sigma}^{*} = e_{\rho + \sigma}^{*}\) ；而最后一个断言可由IV，第50页，公式(12)得出。

注1. --- 同样地，我们可以看到，若 M 是有限生成的自由 A-模，则III，第593页所定义的分次代数 \(\mathbf{S}(\mathbf{M})^{*\mathrm{gr}}\) 可与 \(\mathbf{TS}(\mathbf{M}^*)\) 等同。

命题21. --- A-模的典范同态（IV，第55页）

\[
u: \mathbf {T S} (\mathrm{M}) ^ {* \mathrm{gr}} \rightarrow \operatorname{Pol} _ {\mathrm{A}} (\mathrm{M}, \mathrm{A})
\]

是一个代数同态。

设 \(a \in TS^{q}(M)^{*}\) , \(b \in TS^{r}(M)^{*}\) , \(x \in M\) ；我们有

\[
\begin{array}{r l} u (a b) (x) & = \langle a b, \gamma_ {q + r} (x) \rangle = \langle a \otimes b, c (\gamma_ {q + r} (x)) \rangle = \\ & = \langle a \otimes b, \gamma_ {q} (x) \otimes \gamma_ {r} (x) \rangle = \langle a, \gamma_ {q} (x) \rangle \langle b, \gamma_ {r} (x) \rangle = u (a) (x). u (b) (x), \end{array}
\]

由此即得结论。

注记。--- 2) 复合同态 \(\lambda_{\mathbf{M}}=u\circ\theta:\mathbf{S}(\mathbf{M}^{*})\to\mathrm{Pol}_{\mathbf{A}}(\mathbf{M},\mathbf{A})\) 是唯一的保单位代数同态，它诱导了

\[
\mathbf {M} ^ {*} = \operatorname{Pol} ^ {1} (\mathbf {M}, \mathbf {A})
\]

在 \(\operatorname{Pol}(\mathbf{M}, \mathbf{A})\) 中的包含映射。若 \(\mathbf{M}\) 是有限生成的自由模，且 \(\mathbf{A}\) 是无限整环，则 \(\lambda_{\mathbf{M}}\) 是双射的（命题20及命题19的推论）。特别地，若 \(\mathbf{A}\) 是无限整环，则典范同态 \(f \mapsto \tilde{f}\) 从 \(\mathbf{A}[\mathbf{X}_1, \ldots, \mathbf{X}_n]\) 到 \(\operatorname{Pol}(\mathbf{A}^n, \mathbf{A})\) （IV，第4页）是同构。

\begin{enumerate}
\def\labelenumi{\arabic{enumi})}
\setcounter{enumi}{2}
\tightlist
\item
  考虑余积 \(c_{\mathbf{S}}: \mathbf{S}(\mathbf{M}^{*}) \to \mathbf{S}(\mathbf{M}^{*} \times \mathbf{M}^{*})\) （III，第575页，例6）。对任意 \(v \in \mathbf{S}(\mathbf{M}^{*})\) , \(x, y \in \mathbf{M}\) ，多项式映射 \(\lambda_{\mathbf{M} \times \mathbf{M}}(c_{\mathbf{S}}(v)): \mathbf{M} \times \mathbf{M} \to \mathbf{A}\) 把 \((x, y)\) 映为 \(\lambda_{\mathbf{M}}(v)(x + y)\) 。如此定义的从 \(\mathbf{S}(\mathbf{M}^{*})\) 到 \(\operatorname{Map}(\mathbf{M} \times \mathbf{M}, \mathbf{A})\) 的两个代数同态在 \(\mathbf{M}^{*}\) 上一致，这依据关系式
\end{enumerate}

\[
(v \otimes 1 + 1 \otimes v) (x, y) = v (x + y) \quad (v \in \mathbf {M} ^ {*}).
\]

\section{对称函数}

\subsection{对称多项式}

设 \(n\) 为正整数。对每个置换 \(\sigma \in \mathfrak{S}_n\) ，令 \(\varphi_{\sigma}\) 为 A-代数 \(\mathbf{A}[\mathbf{X}_1, \ldots, \mathbf{X}_n]\) 的这样一个自同构：它把 \(\mathbf{X}_i\) 映为 \(\mathbf{X}_{\sigma(i)}\) ， \(1 \leqslant i \leqslant n\) 。显然 \(\sigma \mapsto \varphi_{\sigma}\) 是从 \(\mathfrak{S}_n\) 到 \(\mathbf{A}[\mathbf{X}_1, \ldots, \mathbf{X}_n]\) 的自同构群中的同态。我们记 \(\sigma f = \varphi_{\sigma}(f)\) ，其中 \(\sigma \in \mathfrak{S}_n\) ， \(f \in \mathbf{A}[\mathbf{X}_1, \ldots, \mathbf{X}_n]\) 。多项式 \(f\) 称为对称的，如果 \(\sigma f = f\) 对所有 \(\sigma \in \mathfrak{S}_n\) 成立；对称多项式构成 \(\mathbf{A}[\mathbf{X}_1, \ldots, \mathbf{X}_n]\) 的一个含单位元的分次子代数；在本段其余部分，我们将其记为 \(\mathbf{A}[\mathbf{X}_1, \ldots, \mathbf{X}_n]^{\text{sym}}\) 。

对每个正整数 \(k\) ，我们用 \(\mathfrak{P}_k\) 表示由 \(k\) 元子集组成的、集合 \(\{1, 2, \ldots, n\}\) 的子集族，并令

\[
s _ {k} = \sum_ {\mathrm{H} \in \mathfrak {P} _ {k}} \prod_ {i \in \mathrm{H}} X _ {i}.\tag{1}
\]

当我们要指明整数 \(n\) 时，我们将用 \(s_{k,n}\) 代替 \(s_k\) 来记。特别地，我们有

\[
\begin{array}{l} s _ {0} = 1 \\ s _ {1} = \sum_ {1 \leqslant i \leqslant n} X _ {i} \\ s _ {2} = \sum_ {1 \leqslant i <   j \leqslant n} X _ {i} X _ {j} \\ \dots \dots \\ s _ {n} = X _ {1} \dots X _ {n} \end{array}
\]

且 \(s_k = 0\) 对 \(k > n\) 成立。显然 \(s_k\) 是 \(k\) 次齐次对称多项式；我们称其为 \(k\) 次初等对称多项式。

在环 \(\mathbf{A}[\mathbf{X}_1,\dots ,\mathbf{X}_n,\mathbf{U},\mathbf{V}]\) 中有关系式

\[
\prod_ {i = 1} ^ {n} \left(\mathrm{U} + \mathrm{VX} _ {i}\right) = \sum_ {k = 0} ^ {n} \mathrm{U} ^ {n - k} \mathrm{V} ^ {k} s _ {k};\tag{2}
\]

通过适当的代入，我们导出关系式

\[
\prod_ {i = 1} ^ {n} \left(1 + \mathrm{TX} _ {i}\right) = \sum_ {k = 0} ^ {n} s _ {k} \mathrm{T} ^ {k},\tag{3}
\]

\[
\prod_ {i = 1} ^ {n} \left(\mathbf {X} - \mathbf {X} _ {i}\right) = \sum_ {k = 0} ^ {n} (- 1) ^ {n - k} s _ {n - k} \mathbf {X} ^ {k}.\tag{4}
\]

定理1. --- 设 \(E = A{[}X_{1}, \ldots, X_{n}{]}\) ， \(S = A{[}X_{1}, \ldots, X_{n}{]}^{\text{sym}}\) 。

\begin{enumerate}
\def\labelenumi{\alph{enumi})}
\item
  对称多项式构成的 A-代数 S 由 \(s_1, \ldots, s_n\) 生成。
\item
  元素 \(s_1, \ldots, s_n\) （属于 \(\mathbf{E}\) ）在 \(\mathbf{A}\) 上代数无关（IV，第4页）。
\item
  单项式族 \(\mathbf{X}^{\nu} = \mathbf{X}_1^{\nu(1)} \ldots \mathbf{X}_n^{\nu(n)}\) （其中 \(0 \leqslant \nu(i) < i\) ， \(1 \leqslant i \leqslant n\) ）是 S-模 \(\mathbf{E}\) 的一组基。特别地， \(\mathbf{E}\) 是秩为 \(n!\) 的自由模。
\end{enumerate}

我们对 \(n\) 用归纳法来证明该定理，情形 \(n = 0\) 是平凡的。记 \(\mathbf{B} = \mathbf{A}[\mathbf{X}_n]\) ，并用 \(s_k'\) 表示 \(k\) 次的、关于 \(\mathbf{X}_1, \ldots, \mathbf{X}_{n-1}\) 的初等对称多项式；于是我们有 \(\mathbf{B}[\mathbf{X}_1, \ldots, \mathbf{X}_{n-1}] = \mathbf{A}[\mathbf{X}_1, \ldots, \mathbf{X}_n]\) 。如果在定理1的陈述中把 \(n\) 换成 \(n-1\) ，把 \(\mathbf{A}\) 换成 \(\mathbf{B}\) ，则归纳假设可表述如下：

\begin{enumerate}
\def\labelenumi{(\Alph{enumi})}
\item
  B-代数 \(S'\) （即由多项式 \(f \in A[X_1, \ldots, X_n]\) 中在 \(X_1, \ldots, X_{n-1}\) 的所有置换下不变者组成）由 \(s_1', \ldots, s_{n-1}'\) 生成。
\item
  元素 \(s_1', \ldots, s_{n-1}'\) （属于 \(E\) ）在 \(B\) 上代数无关。
\item
  单项式族 \(\mathbf{X}_{1}^{\nu(1)}\ldots\mathbf{X}_{n-1}^{\nu(n-1)}\) （其中 \(0\leqslant\nu(i)<i\) ， \(1\leqslant i\leqslant n-1\) ）是 S'-模 E 的一组基。
\end{enumerate}

我们有显然的关系

\[
s _ {k} = s _ {k} ^ {\prime} + s _ {k - 1} ^ {\prime} \mathbf {X} _ {n} \quad (1 \leqslant k \leqslant n - 1),\tag{5}
\]

由此，对 \(k\) 用归纳法得到

\[
s _ {k} ^ {\prime} = (- 1) ^ {k} \mathbf {X} _ {n} ^ {k} + \sum_ {i = 1} ^ {k} (- 1) ^ {k - i} s _ {i} \mathbf {X} _ {n} ^ {k - i} \quad (1 \leqslant k \leqslant n - 1).\tag{6}
\]

我们有 \(\mathbf{S} \subset \mathbf{S}'\) ，因此 \(s_1, \ldots, s_n\) 属于 \(\mathbf{S}'\) ；由(A)和公式(6)，B-代数 \(\mathbf{S}'\) 由 \(s_1, \ldots, s_{n-1}\) 生成。

由(B)，存在一个自同态 \(u\) ，它定义在 B-代数 \(S'\) 上，使得

\[
u \left(s _ {k} ^ {\prime}\right) = (- 1) ^ {k} X _ {n} ^ {k} + \sum_ {i = 1} ^ {k} (- 1) ^ {k - i} s _ {i} ^ {\prime} X _ {n} ^ {k - i} \quad (1 \leqslant k \leqslant n - 1).\tag{7}
\]

根据(5)，我们有 \(u(s_k) = u(s_k') + u(s_{k-1}') \mathbf{X}_n\) ，经简单计算即得 \(u(s_k) = s_k'\) 。设 \(\mathbf{P} \in \mathbf{B}[\mathbf{Y}_1, \ldots, \mathbf{Y}_{n-1}]\) ；于是从 \(\mathbf{P}(s_1, \ldots, s_{n-1}) = 0\) 可推出

\[
0 = u \left(\mathrm{P} \left(s _ {1}, \dots , s _ {n - 1}\right)\right) = \mathrm{P} \left(s _ {1} ^ {\prime}, \dots , s _ {n - 1} ^ {\prime}\right),
\]

因此由(B)得 P = 0。由此可知，B-代数 \(S'\) 由代数无关的元素 \(s_{1}, \ldots, s_{n-1}\) 生成。这一性质可以重新表述如下：

\begin{enumerate}
\def\labelenumi{(\Alph{enumi})}
\setcounter{enumi}{3}
\tightlist
\item
  A-代数 \(S'\) 由代数无关的元素 \(s_1, \ldots, s_{n-1}, X_n\) 生成。
\end{enumerate}

因此我们可以把 \(S'\) 与多项式环 \(C[X_n]\) 等同起来，其中 \(C\) 是 \(E\) 的、由 \(s_1, \ldots, s_{n-1}\) 生成的子 A-代数。

为证明 a)，设 \(f \in S\) 是 \(m\) 次齐次对称多项式。我们有 \(f \in S' = C[X_n]\) ，因此存在一个元素 \(g = P(s_1, \ldots, s_{n-1})\) （属于 \(C\) ），它是 \(m\) 次齐次的（关于 \(X_1, \ldots, X_n\) ），使得 \(f - g\) 可被 \(X_n\) 整除。由于 \(f - g\) 是对称的，它的每一项也可被 \(X_1, \ldots, X_{n-1}\) 整除，因此 \(f - g\) 可被 \(s_n = X_1 \ldots X_n\) 整除。换言之，存在 \(h \in S\) 使得 \(f = g + hs_n\) ，从而 \(\deg h < m\) 。通过对 \(m\) 作归纳即得 \(f\) 属于

\[
\mathbf {C} \left[ s _ {n} \right] _ {\mathrm{E}} = \mathbf {A} \left[ s _ {1}, \dots , s _ {n - 1}, s _ {n} \right] _ {\mathrm{E}}.
\]

这样我们已证明 A-代数 S 由元素 \(s_{1}, \ldots, s_{n}\) 生成。

接下来我们证明 \(b\) ). 若在(4)中以 \(\mathbf{X}_n\) 代替 \(\mathbf{X}\) ，我们得到

\[
(- 1) ^ {n + 1} s _ {n} = \mathrm{X} _ {n} ^ {n} + \sum_ {k = 1} ^ {n - 1} (- 1) ^ {n - k} s _ {n - k} \mathrm{X} _ {n} ^ {k};
\]

换言之， \((-1)^{n + 1}s_n\) 是 \(X_{n}\) 的一个次数为 \(n\) 、系数在 C 中的首一多项式。由IV，第11页，我们于是有以下性质：

\begin{enumerate}
\def\labelenumi{(\Alph{enumi})}
\setcounter{enumi}{4}
\tightlist
\item
  C-代数同态 \(\varphi\) ：从 C{[}T{]} 到 C{[}X \(_n\) {]} = S'，它把 T 映到 \(s_n\) ，是单射，且 S' 是 \(\varphi\) 的像上的自由模，以 (1, X \(_n\) , \ldots, X \(_n^{n-1}\) ) 为基。
\end{enumerate}

元素 \(s_{1}, \ldots, s_{n-1}\) 在 A 上代数无关，由 (D) 可知；因此 \(\varphi\) 的单射性意味着 \(s_{1}, \ldots, s_{n-1}, s_{n}\) 在 A 上代数无关，从而得出 b)。

为证明 c)， \(\varphi\) 的像等于 \(C[s_{n}]_{E} = S\) ，因此由 (E) 可知， \(S'\) 是 S 上以 \((1, X_{n}, ..., X_{n}^{n-1})\) 为基的自由模。现在结论 c) 由归纳假设 (C) 及 II 第 222 页命题 25 得出。证明完毕。

设 \(f\) 是 \(\mathbf{X}_1, \ldots, \mathbf{X}_n\) 中的 \(m\) 次齐次对称多项式。由定理 1（IV，第 62 页），存在多项式 \(\mathbf{Q} \in \mathbf{A}[\mathbf{Y}_1, \ldots, \mathbf{Y}_n]\) 使得 \(f = \mathbf{Q}(s_1, \ldots, s_n)\) 。前面的证明提供了一种显式计算 \(\mathbf{Q}\) 的方法，即对 \(n\) 与 \(m\) 作双重归纳。因为正如我们所见，存在多项式 \(\mathbf{P} \in \mathbf{A}[\mathbf{Y}_1, \ldots, \mathbf{Y}_{n-1}]\) 以及多项式 \(h\) ，它关于 \(\mathbf{X}_1, \ldots, \mathbf{X}_n\) 对称且为 \(m - n\) 次齐次，使得

\[
f = \mathrm{P} (s _ {1}, \dots , s _ {n - 1}) + s _ {n} h.\tag{8}
\]

对于每个多项式 \(u \in \mathbf{A}[\mathbf{X}_1, \ldots, \mathbf{X}_n]\) 我们令

\[
u ^ {\prime} \left(\mathrm{X} _ {1}, \dots , \mathrm{X} _ {n - 1}\right) = u \left(\mathrm{X} _ {1}, \dots , \mathrm{X} _ {n - 1}, 0\right).
\]

则 \(s_1', \ldots, s_{n-1}'\) 是 \(X_1, \ldots, X_{n-1}\) 中的初等对称多项式，且公式(8)蕴含

\[
f ^ {\prime} = \mathrm{P} \left(s _ {1} ^ {\prime}, \dots , s _ {n - 1} ^ {\prime}\right).
\]

因此，P的确定归结为对n - 1个未定元中的对称多项式的计算，而h由(8)得到。

我们通过两个例子来说明该方法。

例。--- 1) 设 \(n = 3\) ，且

\[
f = \mathbf {X} _ {1} ^ {2} (\mathbf {X} _ {2} + \mathbf {X} _ {3}) + \mathbf {X} _ {2} ^ {2} (\mathbf {X} _ {3} + \mathbf {X} _ {1}) + \mathbf {X} _ {3} ^ {2} (\mathbf {X} _ {1} + \mathbf {X} _ {2}).
\]

我们有

\[
f ^ {\prime} = \mathbf {X} _ {1} ^ {2} \mathbf {X} _ {2} + \mathbf {X} _ {1} \mathbf {X} _ {2} ^ {2} = \mathbf {X} _ {1} \mathbf {X} _ {2} (\mathbf {X} _ {1} + \mathbf {X} _ {2}) = s _ {1} ^ {\prime} s _ {2} ^ {\prime}.
\]

我们写 \(g = f - s_1s_2\) ；我们有

\[
\boldsymbol {g} = f - \left(\mathbf {X} _ {1} + \mathbf {X} _ {2} + \mathbf {X} _ {3}\right) \left(\mathbf {X} _ {1} \mathbf {X} _ {2} + \mathbf {X} _ {1} \mathbf {X} _ {3} + \mathbf {X} _ {2} \mathbf {X} _ {3}\right) = - 3 \mathbf {X} _ {1} \mathbf {X} _ {2} \mathbf {X} _ {3},
\]

因此最终

\[
f = s _ {1} s _ {2} - 3 s _ {3}.
\]

\begin{enumerate}
\def\labelenumi{\arabic{enumi})}
\setcounter{enumi}{1}
\tightlist
\item
  再次令 \(n = 3\) 并设 \(p = \mathbf{X}_1^3 + \mathbf{X}_2^3 + \mathbf{X}_3^3\) .
\end{enumerate}

我们有 \(p(\mathbf{X}_1, 0, 0) = \mathbf{X}_1^3 = s_1(\mathbf{X}_1, 0, 0)^3\) . 写出 \(q = p - s_1^3\) ，我们得到

\[
q = - 3 f - 6 \mathrm{X} _ {1} \mathrm{X} _ {2} \mathrm{X} _ {3} = - 3 s _ {1} s _ {2} + 3 s _ {3}
\]

并且最终

\[
p = s _ {1} ^ {3} - 3 s _ {1} s _ {2} + 3 s _ {3}.
\]

设 \(S_1, \ldots, S_n\) 为不定元。我们将给多项式代数 \(A[S_1, \ldots, S_n]\) 配备 N 型分次，其中 \(S_k\) 的权为 k （ \(1 \leqslant k \leqslant n\) ）（IV，第3页），并给 \(A[X_1, \ldots, X_n]\) 配备通常的分次。对 \(1 \leqslant k \leqslant n\) ，初等对称多项式 \(s_{k,n}\) （在 \(X_1, \ldots, X_n\) 中）是 k 次齐次的。根据定理1（IV，第62页），映射 \(g \mapsto g(s_{1,n}, \ldots, s_{n,n})\) 因而是分次代数的同构

\[
\varphi_ {n}: \mathrm{A} [ \mathrm{S} _ {1}, \dots , \mathrm{S} _ {n} ] \rightarrow \mathrm{A} [ \mathrm{X} _ {1}, \dots , \mathrm{X} _ {n} ] ^ {\text { sym }}.
\]

设 \(m\) 为一个满足以下条件的整数 \(0 \leqslant m \leqslant n\) 。对于每个整数 \(k\) 使得 \(1 \leqslant k \leqslant m\) 我们有

\[
s _ {k, m} (\mathbf {X} _ {1}, \dots , \mathbf {X} _ {m}) = s _ {k, n} (\mathbf {X} _ {1}, \dots , \mathbf {X} _ {m}, 0, \dots , 0)\tag{9}
\]

这由 \(s_k\) 的定义（IV，第61页，公式（1））得出。因此图表

\[
\begin{array}{c} \mathrm{A} [ \mathrm{S} _ {1}, \dots , \mathrm{S} _ {m} ] \xrightarrow {j} \mathrm{A} [ \mathrm{S} _ {1}, \dots , \mathrm{S} _ {n} ] \\ \varphi_ {m} \Bigg | _ {\downarrow} \qquad \qquad \qquad \qquad \qquad \qquad \qquad \qquad \qquad \qquad \varphi_ {n} \\ \mathrm{A} [ \mathrm{X} _ {1}, \dots , \mathrm{X} _ {m} ] ^ {\text { sym }} \xleftarrow {\rho} \mathrm{A} [ \mathrm{X} _ {1}, \dots , \mathrm{X} _ {n} ] ^ {\text { sym }} \end{array}\tag{10}
\]

（其中 \(j\) 表示典范包含， \(\rho\) 表示同态

\[
g \mapsto g \left(\mathrm{X} _ {1}, \dots , \mathrm{X} _ {m}, 0, \dots , 0\right))
\]

是交换的。

命题1. --- 对每一对正整数 \(k\) , \(n\) ，令 \(\mathbf{S}_{k}^{(n)}\) 为由 \(\mathbf{X}_{1},\ldots,\mathbf{X}_{n}\) 中所有 \(k\) 次齐次对称多项式组成的 A-模。若整数 \(m\) 满足 \(0\leqslant k\leqslant m\leqslant n\) ，则映射 \(f\mapsto f(\mathbf{X}_{1},\ldots,\mathbf{X}_{m},0,\ldots,0)\) 是 \(\mathbf{S}_{k}^{(n)}\) 到 \(\mathbf{S}_{k}^{(m)}\) 上的同构。

由图表（10）的交换性可知，只需证明权为 k 的、 \(S_{1}, \ldots, S_{n}\) 中的每个等权多项式仅依赖于 \(S_{1}, \ldots, S_{m}\) ，这里假设 \(0 \leqslant k \leqslant m \leqslant n\) 。现在，单项式 \(\mathrm{S}_{1}^{\alpha(1)} \ldots \mathrm{S}_{n}^{\alpha(n)}\) 的权等于整数 \(\alpha(1) + 2\alpha(2) + \cdots + n\alpha(n)\) ；由于整数 \(\alpha(1), \ldots, \alpha(n)\) 为正，关系

\[
\alpha (1) + 2 \alpha (2) + \dots + n \alpha (n) = k \leqslant n
\]

蕴含 \(\alpha(j) = 0\) 对 \(k < j \leqslant n\) 成立，由此得结论。

例3.--- 由IV第64页例2及上述命题1，我们有

\[
\sum_ {i = 1} ^ {n} X _ {i} ^ {3} = s _ {1, n} ^ {3} - 3 s _ {1, n} s _ {2, n} + 3 s _ {3, n}
\]

对于每个整数 \(n \geqslant 3\) 。图(10)的交换性还给出公式

\[
\begin{array}{c} \mathbf {X} _ {1} ^ {3} + \mathbf {X} _ {2} ^ {3} = s _ {1, 2} ^ {3} - 3 s _ {1, 2} s _ {2, 2}, \\ \mathbf {X} _ {1} ^ {3} = s _ {1, 1} ^ {3}. \end{array}
\]

注：--- 设 n 和 k 是两个正整数。我们用 \(\Delta_{k,n}\) 表示 \(N^{n}\) 中长度为 k 的元素组成的集合，并进一步给 \(\Delta_{k,n}\) 配备序关系，记作 \(\alpha \preccurlyeq \beta\) ，它由 \(N^{n}\) 上的字典序（集合论，III，第157页）诱导，并定义群 \(S_{n}\) 在 \(N^{n}\) 上的作用： \((\sigma\alpha)(i) = \alpha(\sigma^{-1}(i))\) ，其中 \(\sigma \in S_{n}\) , \(\alpha \in N^{n}\) ， \(1 \leqslant i \leqslant n\) 。此外，我们用 \(D_{k}\) 表示满足下述条件的元素 \(\alpha = (\alpha(1), \ldots, \alpha(k))\) （属于 \(N^{k}\) ）组成的集合：

\[
\alpha (1) \geqslant \alpha (2) \geqslant \dots \geqslant \alpha (k), \quad \alpha (1) + \dots + \alpha (k) = k.
\]

假设 \(k \leqslant n\) ，并将 \(\mathbf{N}^k\) 与 \(\mathbf{N}^n\) 的一个子集经由映射 \((\alpha(1), \ldots, \alpha(k)) \mapsto (\alpha(1), \ldots, \alpha(k), 0, \ldots, 0)\) 等同起来。则 \(D_k\) 由元素 \(\alpha\) （属于 \(\Delta_{k,n}\) ）组成，使得 \(\sigma \alpha \preccurlyeq \alpha\) 对所有 \(\sigma \in \mathfrak{S}_n\) 成立。由此可知，群 \(\mathfrak{S}_n\) 在 \(\Delta_{k,n}\) 中的每个轨道都包含 \(D_k\) 的唯一一个元素。对每个 \(\alpha \in D_k\) ，令 \(O(\alpha)\) 为 \(\alpha\) 在 \(\Delta_{k,n}\) 中在 \(\mathfrak{S}_n\) 作用下的轨道；令

\[
\mathbf {M} (\alpha) = \sum_ {\beta \in O (\alpha)} \mathbf {X} ^ {\beta}.\tag{11}
\]

由IV第47页引理1可知，族 \((\mathbf{M}(\alpha))_{\alpha \in \mathrm{D}_k}\) 是 A-模 \(\mathbf{S}_k^{(n)}\) 的一组基。

对每个 \(\alpha \in D_k\) ，令

\[
\mathrm{S} (\alpha) = \prod_ {i = 1} ^ {k} s _ {i} ^ {\alpha (i) - \alpha (i + 1)} \quad (\text { where   it   is   understood   that } \alpha (k + 1) = 0);\tag{12}
\]

由于我们有 \(\sum_{i=1}^{k} i \cdot (\alpha(i) - \alpha(i+1)) = \sum_{i=1}^{k} \alpha(i) = k\) ，对称多项式 \(S(\alpha)\) 是 \(k\) 次齐次的。由定理1（IV，第62页）直接可得，族 \((S(\alpha))_{\alpha \in D_k}\) 是 A-模 \(S_k^{(n)}\) 的一组基。

设 \(\alpha, \beta\) 属于 \(D_k\) ，并设 \(c_{\alpha\beta}\) 为单项式 \(X^\beta\) 在多项式 \(S(\alpha)\) （在情形 \(A = Z\) 且 \(k = n\) 下由(12)定义）中的系数。它是一个正整数，与环 \(A\) 及整数 \(n\) 无关。由公式(9)（IV，第64页），我们有

\[
\mathrm{S} (\alpha) = \sum_ {\beta \in \mathrm{D} _ {k}} c _ {\alpha \beta} \cdot \mathrm{M} (\beta) \quad (\alpha \in \mathrm{D} _ {k}).\tag{13}
\]

可以证明（参见IV，第101页，练习13），矩阵 \(\mathbf{C} = (c_{\alpha \beta})_{\alpha, \beta \in D_k}\) 满足 \(c_{\alpha \alpha} = 1\) ，且 \(c_{\alpha \beta} = 0\) （对 \(\alpha < \beta\) ）；推广II，第351页中引入的术语，我们称 \(\mathbf{C}\) 属于下全三角群。矩阵 \(\mathbf{D}\) （ \(\mathbf{C}\) 的逆矩阵）也是如此。矩阵 \(\mathbf{C}\) 与 \(\mathbf{D}\) 当 \(2 \leqslant k \leqslant 5\) 时的值列于表（IV，第103-105页）中。

现在假设 \(n < k\) ，并将 \(\mathbf{N}^n\) 与 \(\mathbf{N}^k\) 的一个子集通过映射 \((\alpha(1), \ldots, \alpha(n)) \mapsto (\alpha(1), \ldots, \alpha(n), 0, \ldots, 0)\) 等同起来。对每个 \(\alpha\) （在 \(D_k \cap N^n\) 中），我们再次用 \(O(\alpha)\) 表示 \(\alpha\) 在 \(\Delta_{k,n}\) 中在 \(\mathfrak{S}_n\) 作用下的轨道，并按（11）定义 \(M(\alpha)\) 。进一步，我们按（12）定义 \(S(\alpha)\) ，则族 \((M(\alpha))_{\alpha \in D_k \cap N^n}\) 与 \((S(\alpha))_{\alpha \in D_k \cap N^n}\) 是 A-模 \(S_k^{(n)}\) 的基。根据 IV 第64页的公式（9），我们有一个与（13）类似的公式，其中 \(D_k\) 被换成 \(D_k \cap N^n\) ，且具有相同的整数 \(c_{\alpha\beta}\) 。

例4. --- 根据 IV 第65页的例3，我们有

\[
\mathbf {M} (3, 0, 0) = \mathbf {S} (3, 0, 0) - 3 \mathbf {S} (2, 1, 0) + 3 \mathbf {S} (1, 1, 1)
\]

对每个整数 \(n \geqslant 3\) 成立，并且

\[
\mathbf {M} (3, 0) = \mathbf {S} (3, 0) - 3 \mathbf {S} (2, 1)
\]

对于 n = 2。

\subsection{对称有理分式}

设 K 为一个交换域，且 \(X_{1}, X_{2}, \ldots, X_{n}\) 为不定元。对每个 \(\sigma \in S_{n}\) ，我们在第 1 号（IV，第61页）中定义了一个自同构 \(\varphi_{\sigma}\) （作用于 \(K[X_{1}, X_{2}, \ldots, X_{n}]\) ）。该自同构具有唯一的扩张 \(\psi_{\sigma}\) （到域 \(K(X_{1}, \ldots, X_{n})\) 上），且 \(\sigma \mapsto \psi_{\sigma}\) 是从 \(S_{n}\) 到 \(K(X_{1}, \ldots, X_{n})\) 的自同构群中的单同态。对每个 \(f \in K(X_{1}, \ldots, X_{n})\) ，我们有 \((\psi_{\sigma} f) (X_{1}, \ldots, X_{n}) = f(X_{\sigma(1)}, \ldots, X_{\sigma(n)})\) 。使得 \(\psi_{\sigma}(f) = f\) （对所有 \(\sigma \in S_{n}\) ）的有理分式 f 称为对称有理分式。所有对称有理分式（在 \(X_{1}, \ldots, X_{n}\) 中）组成的集合是 \(K(X_{1}, \ldots, X_{n})\) 的一个子域。

命题2. --- \(X_{1}, \ldots, X_{n}\) 中对称有理分式组成的域是 \(X_{1}, \ldots, X_{n}\) 中对称多项式环的分式域。

设 \(f \in \mathbf{K}(\mathbf{X}_1, \ldots, \mathbf{X}_n)\) 是一个对称有理分式，且设 \(u_1, v_1\) 是 \(\mathbf{K}[\mathbf{X}_1, \ldots, \mathbf{X}_n]\) 中使得 \(f = \frac{u_1}{v_1}\) 的两个元素。我们令 \(v = \prod_{\sigma \in \mathfrak{S}_n} \psi_\sigma(v_1) \in \mathbf{K}[\mathbf{X}_1, \ldots, \mathbf{X}_n]\) ， \(u = fv \in \mathbf{K}[\mathbf{X}_1, \ldots, \mathbf{X}_n]\) ；则 \(v\) 是对称的，因而 \(u\) 也是对称的（因为 \(f\) 是），且 \(f = \frac{u}{v}\) ，由此即得结论。

推论。--- 设 \(s_{1}, s_{2}, \ldots, s_{n}\) 是 \(X_{1}, \ldots, X_{n}\) 中的初等对称多项式。对每个有理分式 \(g \in \mathbf{K}(S_{1}, S_{2}, \ldots, S_{n})\) ，序列 \((s_{1}, s_{2}, \ldots, s_{n})\) 在 g 中可代入，且映射 \(g \mapsto g(s_{1}, s_{2}, \ldots, s_{n})\) 是 \(\mathbf{K}(S_{1}, S_{2}, \ldots, S_{n})\) 到 \(X_{1}, \ldots, X_{n}\) 中对称有理分式域上的同构。

这由IV第62页的命题2和定理1得出。

\subsection{对称形式幂级数}

设 I 为一个集合， \(\mathbf{X} = (\mathbf{X}_i)_{i\in I}\) 为一族未定元， \(\mathbf{A}[[\mathbf{X}]]\) 为 \(\mathbf{X}_i\) 中的形式幂级数代数。对每个置换 \(\sigma \in \mathfrak{S}_I\) ，存在唯一的连续自同构 \(\varphi_{\sigma}\) （作用于代数 \(\mathbf{A}[[\mathbf{X}]]\) ），它把每个 \(\mathbf{X}_i\) 映为 \(\mathbf{X}_{\sigma (i)}\) （ \(i\in I\) ）（IV，第28页，命题4）；显然 \(\sigma \mapsto \varphi_{\sigma}\) 是从 \(\mathfrak{S}_I\) 到代数 \(\mathbf{A}[[\mathbf{X}]]\) 的连续自同构群中的同态。设 \(f\in \mathbf{A}[[\mathbf{X}]]\) 为一个形式幂级数；我们令 \(\sigma f = \varphi_{\sigma}(f)\) ，并称形式幂级数 \(f\) 是对称的，如果 \(\sigma f = f\) （对每个 \(\sigma \in \mathfrak{S}_I\) ）。对称形式幂级数构成 \(\mathbf{A}[[\mathbf{X}]]\) 的一个闭子代数，记为 \(\mathbf{A}[[\mathbf{X}]]^{\mathrm{sym}}\) ，并配备由 \(\mathbf{A}[[\mathbf{X}]]\)

诱导的拓扑。设 T 为一个未定元。在形式幂级数环 A{[}{[}X, T{]}{]} 中，族 \((\mathbf{X}_i\mathbf{T})_{i\in I}\) 是可求和的，因而族 \((1 + \mathbf{X}_i\mathbf{T})_{i\in I}\) 是可乘的（IV，第27页，命题2）；此外我们有

\[
\prod_ {i \in I} (1 + X _ {i} T) = 1 + \sum_ {k \geqslant 1} s _ {k} T ^ {k},\tag{14}
\]

其中形式幂级数 \(s_k \in \mathbf{A}[[\mathbf{X}]]\) 由下式定义

\[
s _ {k} = \sum_ {\mathrm{H} \in \mathfrak {P} _ {k}} \left(\prod_ {i \in \mathrm{H}} \mathrm{X} _ {i}\right) \quad (k \geqslant 1)\tag{15}
\]

（我们用 \(\mathfrak{P}_k\) 表示 I 的所有 \(k\) 元子集组成的集合）。特别地，我们有 \(s_1 = \sum_{i \in I} X_i\) 。当 I 是含 \(n\) 个元素的有限集时，我们有 \(s_k = 0\) 对 \(k > n\) 成立；更精确地说，当 \(\mathbf{I} = \{1,\dots,n\}\) 时，形式幂级数 \(s_k\) 正是 \(k\) 次的初等对称多项式（在 \(\mathbf{X}_1,\ldots,\mathbf{X}_n\) 中）。

设 \(\mathbf{S} = (\mathbf{S}_k)_{k\geqslant 1}\) 为一个未定元序列。由于形式幂级数 \(s_k\) 的阶 \(\geqslant k\) ，且属于 \(\mathbf{A}[[\mathbf{X}]]^{\mathrm{sym}}\) ，IV第28页命题4的条件 a) 和 b) 在 \(\mathrm{E} = \mathbf{A}[[\mathbf{X}]]^{\mathrm{sym}}\) 下得到满足；因此存在唯一的连续 A-代数同态

\[
\varphi_ {I}: \mathrm{A} [ [ \mathrm{S} ] ] \rightarrow \mathrm{A} [ [ \mathrm{X} ] ] ^ {\text { sym }}
\]

使得 \(\varphi_{I}(S_{k}) = s_{k}\) 对每个整数 \(k\geqslant 1\) 成立。

定理2. --- a) 若 I 是含 n 个元素的有限集，则 \(\varphi_{I}\) 诱导出从 A{[}{[}S \(_{1}\) , \ldots，S \(_{n}\) {]}{]} 到 A{[}{[}X{]}{]} \(^{\text{sym}}\) 上的双连续同构。

\begin{enumerate}
\def\labelenumi{\alph{enumi})}
\setcounter{enumi}{1}
\tightlist
\item
  如果 I 是无限集， \(\varphi_{\mathrm{I}}\) 是从 A{[}{[}S{]}{]} 到 A{[}{[}X{]}{]} \(^{\mathrm{sym}}\) 上的双连续同构。
\end{enumerate}

在情形 a) 中，我们令 \(\mathbf{B} = \mathbf{A}[[\mathbf{S}_1, \ldots, \mathbf{S}_n]]\) ，并赋予该代数由 \(\mathbf{A}[[\mathbf{S}]]\) 的拓扑诱导的拓扑；我们还用 \(\psi_{\mathrm{I}}\) 表示 \(\varphi_{\mathrm{I}}\) 在 \(\mathbf{B}\) 上的限制。在情形 b) 中，我们令 \(\mathbf{B} = \mathbf{A}[[\mathbf{S}]]\) ，且 \(\psi_{\mathrm{I}} = \varphi_{\mathrm{I}}\) 。我们将给多项式代数 \(\mathbf{A}[\mathbf{S}]\) 配备 \(\mathbf{N}\) 型分次，使得 \(\mathbf{S}_k\) 的权为 \(k\) （对每个整数 \(k \geqslant 1\) ）。

引理 1. --- 设 J 是 I 的有限子集，r 是使 Card \(J \geqslant r\) 成立的整数，且 f 是关于 \(X_{i}\) ( \(i \in I\) ) 的 r 次齐次对称形式幂级数。设 \(\bar{f}\) 为对 I - J 中每个 i 把 0 代入 \(X_{i}\) 所得的形式幂级数。如果 \(\bar{f} = 0\) ，则我们有 f = 0。

我们令 \(f = \sum_{|\alpha| = r} a_{\alpha} X^{\alpha}\) （其中 \(|\alpha|\) 是长度 \(\sum_{i \in I} \alpha_i\) ，即多重指标 \(\alpha = (\alpha_i)_{i \in I}\) 的长度）。如果 \(\alpha\) 是长度为 \(r\) 的多重指标，且 \(J'\) 是 \(\alpha\) 的支集（即由这样的 \(i \in I\) 组成的集合： \(\alpha_i \neq 0\) ），则 Card \(J' \leqslant r\) ，因此存在一个置换 \(\sigma \in \mathfrak{S}_I\) 使得 \(\sigma(J') \subset J\) 。令 \(\beta_i = \alpha_{\sigma^{-1}(i)}\) （对 \(i \in I\) ），那么单项式 \(X^\beta = \prod_{i \in I} X_{\sigma(i)}^{\alpha_i}\) 仅依赖于不定元 \(X_j (j \in J)\) ，从而 \(a_\beta = 0\) （由假设 \(\bar{f}=0\) ）。由于 f 是对称的，我们有 \(a_{\alpha}=a_{\beta}\) ，且由于 \(\alpha\) 是任意的，我们得出结论 f=0。

引理2. --- 设 \(f\) 是次数为 \(r\) 的、关于 \(X_i (i \in I)\) 的对称形式幂级数。存在唯一的多项式 \(P \in B \cap A[S]\) ，它是权为 \(r\) 的等权多项式，使得 \(f = \psi_I(P)\) 。

当I为有限时，该情形由定理1（IV，第62页）得出。

假设 I 是无限集，并选取 I 的一个含 r 个元素的有限子集 J。我们沿用引理 1 的记号；我们注意到， \(S_{n}\) ( \(n \geqslant 1\) ) 中每个权为 r 的等权多项式仅依赖于 \(S_{1}, \ldots, S_{r}\) ，且 \(\bar{s}_{1}, \ldots, \bar{s}_{r}\) 是 r 个不定元 \(X_{j}\) ( \(j \in J\) ) 中的初等对称多项式。若 P 是 \(S_{n}\) 中权为 r 的等权多项式，且 \(h = f - \psi_{\mathrm{I}}(\mathrm{P})\) ，则 \(\bar{h} = \bar{f} - \mathrm{P}(\bar{s}_{1}, \ldots, \bar{s}_{r})\) ，且引理 1 表明关系 \(f = \psi_{\mathrm{I}}(\mathrm{P})\) 等价于 \(\bar{f} = \mathrm{P}(\bar{s}_{1}, \ldots, \bar{s}_{r})\) 。由定理 1（IV，第 62 页），存在唯一的多项式 \(P \in A[S]\) ，等权且权为 r ，使得 \(\bar{f} = \mathrm{P}(\bar{s}_{1}, \ldots, \bar{s}_{r})\) ，由此即得结论。

引理 3. --- 对每个整数 \(m \geqslant 0\) ，令 \(\mathfrak{c}_m\) 为代数 \(\mathrm{A}[[\mathrm{X}]]^{\mathrm{sym}}\) 的理想，它由所有阶 \(\geqslant m\) 的对称形式幂级数组成。序列 \((\mathfrak{c}_m)_{m \geqslant 0}\) 是 \(\mathrm{A}[[\mathrm{X}]]^{\mathrm{sym}}\) 中 0 的邻域基。

当 I 为有限集时，引理显然成立，因此我们假设 I 为无限集。对 I 的每个由 m 个元素组成的有限子集 J，记 \(\tilde{J}\) 为 \(\mathbf{N}^{(1)}\) 中长度 \textless{} m 且支集含于 J 的元素组成的集合。进一步记 \(a_{j}^{\prime}\) 为不含任何形如 \(aX^{\alpha}\) （其中齐次 \(\alpha \in \tilde{J}\) ）的项的形式幂级数组成的集合。由于 \(\tilde{J}\) 是 \(\mathbf{N}^{(1)}\) 的有限子集，且 \(\mathbf{N}^{(1)}\) 的每个有限子集都包含于形如 \(\tilde{J}\) 的集合，族 \((\mathfrak{a}_{\mathrm{j}}^{\prime})\) 是 A{[}{[}X{]}{]} 中 0 的一个邻域基（IV，第26页）。现在，引理1蕴含关系式 \(a_{j}^{\prime} \cap A[[X]]^{\text{sym}} = c_{m}\) 对每个含 m 个元素的子集 J 成立，这便证明了引理3。

由于在 \(S_{k}\) 中给定权的单项式只有有限多个，每个形式幂级数 \(f \in B\) 都可以唯一地写成 \(f = \sum_{r \geqslant 0} P_{r}\) ，其中 \(P_{r}\) 是 \(B \cap A[S]\) 中权为 r 的等权多项式。对每个整数 \(m \geqslant 0\) ，令 \(\mathfrak{b}_m\) 为 B 中由所有上述类型的、满足 \(P_r = 0\) （对 \(0 \leqslant r < m\) ）的形式幂级数组成的理想。序列 \((\mathfrak{b}_m)_{m \geqslant 0}\) 是 B 中 0 的邻域基。

沿用上述记号， \(\psi_{\mathrm{I}}(\mathbf{P}_{r})\) 是关于 \(\mathbf{X}_i\) 的对称形式幂级数，且是 \(r\) 次齐次的，因而它是 \(r\) 次齐次分量（属于 \(\psi_{\mathrm{I}}(f)\) ）。引理2表明 \(\psi_{\mathrm{I}}\) 是 B 到 \(\mathbf{A}[[\mathbf{X}]]^{\mathrm{sym}}\) 上的代数同态，它把 \(\mathfrak{b}_m\) 映到 \(\mathfrak{c}_m\) （对每个整数 \(m \geqslant 0\) ）；于是引理3表明 \(\psi_{\mathrm{I}}\) 是双连续的，这就完成了定理2的证明。

\subsection{幂的和}

我们再次记 \(\mathbf{X} = (\mathbf{X}_{i})_{i \in I}\) ，表示一族未定元。对称形式幂级数 \(s_{k}\) 如前所定义，

\[
s _ {k} = \sum_ {\mathrm{H} \in \mathfrak {P} _ {k}} \prod_ {i \in \mathrm{H}} X _ {i} \quad (k \geqslant 1),\tag{16}
\]

，其中 \(P_{k}\) 是 I 的所有 k 元子集组成的集合。我们也记

\[
p _ {k} = \sum_ {i \in I} X _ {i} ^ {k} \quad (k \geqslant 1).\tag{17}
\]

这是一个对称形式幂级数，且是 k 次齐次的。

引理4（“牛顿关系”）。--- 对每个整数 \(d \geqslant 1\) ，我们有

\[
p _ {d} = \sum_ {k = 1} ^ {d - 1} (- 1) ^ {k - 1} s _ {k} p _ {d - k} + (- 1) ^ {d + 1} d s _ {d}.\tag{18}
\]

我们定义连续导子 \(\Delta\) （作用于 \(\mathbf{A}[[\mathbf{X}]]\) ）如下： \(\Delta(u) = \sum_{n \geq 0} nu_n\) ，其中对每个 \(u\) （属于 \(\mathbf{A}[[\mathbf{X}]]\) ）， \(u_n\) 表示次数为 \(n\) 的齐次分量（属于 \(u\) ）。由(16)及IV第27页命题2，我们有

\[
1 + \sum_ {k \geqslant 1} s _ {k} = \prod_ {i \in I} (1 + X _ {i}).\tag{19}
\]

于是由IV第34页的命题9，我们有

\[
\left(\sum_ {k \geqslant 1} k s _ {k}\right) \cdot \left(1 + \sum_ {k \geqslant 1} s _ {k}\right) ^ {- 1} = \sum_ {i \in I} \Delta \left(\mathrm{X} _ {i}\right) / \left(1 + \mathrm{X} _ {i}\right).\tag{20}
\]

我们有 \(\Delta(X_i) = X_i\) ，且 \(X_i / (1 + X_i) = \sum_{k \geq 1} (-1)^{k-1} X_i^k\) 。因此(20)的右端等于 \(\sum_{k \geq 1} (-1)^{k-1} p_k\) 。于是由(20)，我们有

\[
\sum_ {k \geqslant 1} k s _ {k} = \left(1 + \sum_ {k \geqslant 1} s _ {k}\right) \cdot \left(\sum_ {k \geqslant 1} (- 1) ^ {k - 1} p _ {k}\right)
\]

于是，通过比较次数为 \(d\) 的齐次分量，即得引理4。

注记。--- 使用上述证明中的记号，我们有

\[
\Delta u = \sum_ {i \in I} X _ {i} \cdot D _ {i} (u)
\]

(IV，第33页，推论1)。换言之，欧拉关系式（IV，第8页，命题6）可以推广到形式幂级数：若 \(u \in \mathbf{A}[\mathbf{X}]\) 是 \(n\) 次齐次的，则

\[
n \cdot u = \sum_ {i \in I} X _ {i} \cdot D _ {i} (u).\tag{21}
\]

当 I 是由 \(n\) 个元素组成的有限集时，我们有 \(s_k = 0\) 对 \(k > n\) 成立。于是牛顿关系式可以写成

\[
\begin{array}{l} p _ {2} = s _ {1} p _ {1} - 2 s _ {2} \\ p _ {3} = s _ {1} p _ {2} - s _ {2} p _ {1} + 3 s _ {3} \\ \dots \dots \\ p _ {n - 1} = s _ {1} p _ {n - 2} - s _ {2} p _ {n - 3} + \dots + (- 1) ^ {n - 1} s _ {n - 2} p _ {1} + (- 1) ^ {n} (n - 1) s _ {n - 1} \\ p _ {n} = s _ {1} p _ {n - 1} - s _ {2} p _ {n - 2} + \dots + (- 1) ^ {n} s _ {n - 1} p _ {1} + (- 1) ^ {n + 1} n s _ {n} \end{array}
\]

以及

\[
p _ {k} = s _ {1} p _ {k - 1} - s _ {2} p _ {k - 2} + \dots + (- 1) ^ {n + 1} s _ {n} p _ {k - n} \quad (\text { for } k > n).\tag{22}
\]

上述关系中的前 n 个对任意 I 均成立；例如我们得到

\[
p _ {1} = s _ {1}, \quad p _ {2} = s _ {1} ^ {2} - 2 s _ {2}, \quad p _ {3} = s _ {1} ^ {3} - 3 s _ {1} s _ {2} + 3 s _ {3}.
\]

更一般地，设 \(\mathbf{S} = (\mathbf{S}_n)_{n\geqslant 1}\) 为一族未定元。我们递归地定义多项式 \(\mathbf{P}_d\in \mathbf{Z}[\mathbf{S}_1,\dots ,\mathbf{S}_d]\) 如下： \(\mathbf{P}_1 = \mathbf{S}_1\) ，且

\[
\mathrm{P} _ {d} = \sum_ {k = 1} ^ {d - 1} (- 1) ^ {k - 1} \mathrm{S} _ {k} \mathrm{P} _ {d - k} + (- 1) ^ {d + 1} d \mathrm{S} _ {d} \quad (d \geqslant 2).
\]

于是我们有以下“普适公式” \(p_d = \mathrm{P}_d(s_1, \ldots, s_d)\) ，它们对任意环 A 及未定元族 X 均成立。

设 \(\mathbf{P} = (\mathrm{P}_k)_{k\geqslant 1}\) 为一列未定元。由于 \(p_k\) 是 \(k\) 次齐次的（在 \(\mathbf{A}[[\mathbf{X}]]\) 中），因此恰好存在一个连续 A-代数同态

\[
\lambda_ {I}: \mathbf {A} [ [ \mathbf {P} ] ] \rightarrow \mathbf {A} [ [ \mathbf {X} ] ] ^ {\text { sym }}
\]

使得 \(\lambda_{I}(\mathbf{P}_{k}) = p_{k}\) （对每个整数 \(k \geqslant 1\) ）（IV，第28页）。若我们赋予 \(\mathbf{P}_{k}\) 权 \(k\) ，则 \(\lambda_{I}\) 把权为 \(n\) 的等权多项式（在 \(\mathbf{P}_{k}\) 中）变换为 \(n\) 次齐次的形式幂级数（在 \(\mathbf{X}_{i}\) 中）。

命题3. --- a) 若 I 是含 n 个元素的有限集，且 \(n! \cdot 1\) 在 A 中可逆，则 \(\lambda_{I}\) 诱导出从 \(A[[P_{1}, ..., P_{n}]]\) 到 \(A[[X]]^{sym}\) 上的双连续同构。

\begin{enumerate}
\def\labelenumi{\alph{enumi})}
\setcounter{enumi}{1}
\tightlist
\item
  如果 I 是无限集且 A 是一个 Q-代数，那么 \(\lambda_{I}\) 是从 A{[}{[}P{]}{]} 到 A{[}{[}X{]}{]} \(^{\text{sym}}\) 上的双连续同构。
\end{enumerate}

我们将仅处理情形 a)，情形 b) 完全类似。

由定理 2（IV，第68页），我们可以把 A{[}{[}X{]}{]} \(^{sym}\) 与形式幂级数代数 A{[}{[}S \(_1\) , \ldots，S \(_n\) {]}{]} 等同起来，其中 S \(_k\) 对应于 s \(_k\) 。根据IV第70页引理4，存在形式幂级数 g \(_1\) , \ldots，g \(_n\) ，其阶 \(\geqslant 2\) ，关于未定元 s \(_1\) , \ldots，s \(_n\) ，使得

\[
p _ {k} = (- 1) ^ {k + 1} k s _ {k} + g _ {k} (s _ {1}, \dots , s _ {n}) \quad (1 \leqslant k \leqslant n).
\]

由于 \(k!\) . 1 在 A 中可逆，IV第35页引理2证明了拓扑 A-代数 \(\mathbf{A}[[\mathbf{X}]]^{\mathrm{sym}}\) 的自同构 T 的存在性，它把 \(s_k\) 映为 \(p_k\) （对 \(1 \leqslant k \leqslant n\) ）。于是命题3 a) 是直接推论。

推论。--- 设 \(\xi_1, \ldots, \xi_n, \eta_1, \ldots, \eta_n\) 是 A 的元素，并假设 A 是整环。

\begin{enumerate}
\def\labelenumi{\alph{enumi})}
\item
  如果 \(s_k(\xi_1, \ldots, \xi_n) = s_k(\eta_1, \ldots, \eta_n)\) 对 \(1 \leqslant k \leqslant n\) 成立，则存在一个置换 \(\sigma \in \mathfrak{S}_n\) ，使得 \(\eta_i = \xi_{\sigma(i)}\) 对 \(1 \leqslant i \leqslant n\) 成立。
\item
  假设 \(n!\) . \(1 \neq 0\) 在 A 中，且
\end{enumerate}

\[
\xi_ {1} ^ {k} + \dots + \xi_ {n} ^ {k} = \eta_ {1} ^ {k} + \dots + \eta_ {n} ^ {k}\tag{23}
\]

对 \(1 \leqslant k \leqslant n\) 成立。则存在一个置换 \(\sigma \in \mathfrak{S}_n\) ，使得 \(\eta_i = \xi_{\sigma(i)}\) 对 \(1 \leqslant i \leqslant n\) 成立。

在假设 \(a)\) 下，我们有 \(\prod_{i=1}^{n} (\mathbf{X} - \xi_i) = \prod_{i=1}^{n} (\mathbf{X} - \eta_i)\) 。如果我们以 \(\eta_n\) 代替 \(\mathbf{X}\) ，则得到 \(\prod_{i=1}^{n} (\eta_n - \xi_i) = 0\) ，且由于 \(\mathbf{A}\) 是整环，存在整数 \(\sigma(n)\) 使得 \(1 \leqslant \sigma(n) \leqslant n\) 且 \(\eta_n = \xi_{\sigma(n)}\) 。于是断言 \(a)\) 由归纳法容易得出，因为 \(\mathbf{A}[\mathbf{X}]\) 是整环。

在假设 \(b)\) 下，由命题3，存在多项式 \(\Pi_1, \ldots, \Pi_n\) （在 \(n\) 个不定元中），其系数属于 \(A\) 的分式域，使得 \(s_k = \Pi_k(p_1, \ldots, p_n)\) 对 \(1 \leqslant k \leqslant n\) 成立。现在关系式(23)蕴含

\[
s _ {k} (\xi_ {1}, \dots , \xi_ {n}) = s _ {k} (\eta_ {1}, \dots , \eta_ {n})
\]

对 \(1 \leqslant k \leqslant n\) 成立，因而 \(b)\) 由 \(a)\) 得出。

\subsection{多项式根中的对称函数}

考虑一个次数为 \(n\) 、系数在 \(A\) 中的首一多项式：

\[
f = \mathbf {X} ^ {n} + a _ {1} \mathbf {X} ^ {n - 1} + \dots + a _ {n - 1} \mathbf {X} + a _ {n}.
\]

我们定义一个结合的、交换的且保单位的 A-代数 \(\mathbf{E}_f\) ，它以 \(x_{1},\ldots ,x_{n}\) 为生成元，关系式为

\[
\sum_ {i _ {1} <   \dots <   i _ {k}} x _ {i _ {1}} \dots x _ {i _ {k}} = (- 1) ^ {k} a _ {k} \quad (1 \leqslant k \leqslant n).\tag{24}
\]

更精确地说，我们有

\[
\mathrm{E} _ {f} = \mathrm{A} [ \mathrm{X} _ {1}, \dots , \mathrm{X} _ {n} ] / \mathfrak {a}
\]

其中理想 \(\mathfrak{a}\) 由多项式 \(s_k + (-1)^{k+1} a_k\) （ \(1 \leqslant k \leqslant n\) ）生成，而 \(x_i\) 是 \(\mathbf{X}_i\) 模 \(\mathfrak{a}\) 的剩余类（ \(1 \leqslant i \leqslant n\) ）。关系式 (24) 也等价于 \(f(\mathbf{X}) = \prod_{i=1}^{n} (\mathbf{X} - x_i)\) 。当有任何歧义风险时，我们写作 \(x_{1,f}, \ldots, x_{n,f}\) 以代替 \(x_1, \ldots, x_n\) 。

命题4. --- 设 B 为一个交换环， \(\rho\) 为从 A 到 B 的一个同态，以及 \(\xi_1, \ldots, \xi_n\) 为 B 的元素。假设关系 \( {}^{\rho}f(X) = \prod_{i=1}^{n}(X - \xi_i)\) 在 B{[}X{]} 中成立。则存在一个且仅一个环同态 \(u: E_f \to B\) ，使得 \(\rho(a) = u(a, 1)\) 对所有 \(a \in A\) 成立，且 \(u(x_i) = \xi_i\) 对 \(1 \leqslant i \leqslant n\) 成立。

我们通过 \(\rho\) 把 B 视为结合的、交换的且保单位的 A-代数。于是关系 \({}^{\rho}f(X) = \prod_{i=1}^{n}(X - \xi_i)\) 也可以写成如下形式

\[
\sum_ {i _ {1} <   \dots <   i _ {k}} \xi_ {i _ {1}} \dots \xi_ {i _ {k}} = (- 1) ^ {k} a _ {k} \cdot 1 \quad (1 \leqslant k \leqslant n)
\]

在 B 中。由于关系式 (24) 给出了 \(\mathbf{E}_f\) 的一个表示，命题4得证。

命题4证明了把 \(E_{f}\) 称为“f 的泛分解代数”这一名称的合理性。关系式 \(f(\mathbf{X}) = \prod_{i=1}^{n} (\mathbf{X} - x_{i,f})\) 称为“f 的泛分解”。设 \(\sigma \in S_{n}\) 为一个置换；由于 \(f(\mathbf{X}) = \prod_{i=1}^{n} (\mathbf{X} - x_{\sigma(i),f})\) ，存在一个自同构 \(t_{\sigma}\) （A-代数 \(E_{f}\) 的），其特征为 \(t_{\sigma}(x_{i,f}) = x_{\sigma(i),f}\) （ \(1 \leqslant i \leqslant n\) ）。我们有 \(t_{\sigma\tau} = t_{\sigma} \circ t_{\tau}\) （对 \(\sigma, \tau\) ，均在 \(S_{n}\) 中），从而得到群 \(S_{n}\) 在 A-代数 \(E_{f}\) 上的作用。

命题5. --- 在泛分解代数 \(E_{f}\) 中，单项式族 \(x_{1}^{\nu(1)}\ldots x_{n}^{\nu(n)}\) （其中 \(0 \leqslant \nu(i) < i\) ， \(1 \leqslant i \leqslant n\) ）是 A-模 \(E_{f}\) 的一组基。特别地， \(E_{f}\) 是秩为 \(n!\)

的自由 A-模。记 \(\mathbf{B} = \mathbf{A}[X_1, \ldots, X_n]\) ，且 \(\mathbf{C} = \mathbf{A}[X_1, \ldots, X_n]^{\mathrm{sym}}\) 。由定理1（IV，第62页），我们有 \(\mathbf{C} = \mathbf{A}[s_1, \ldots, s_n]\) ，且 \(s_1, \ldots, s_n\) 在 \(\mathbf{A}\) 上代数无关。 \(s_1, \ldots, s_n\) 中无常数项的多项式构成一个理想 \(\mathbf{C}^+\) （在 \(\mathbf{C}\) 中），它是 \(\mathbf{A}\) 的补，且由 \(s_1, \ldots, s_n\) 生成。设 \(\mathfrak{c}\) 为 \(\mathbf{C}\) 的由 \(s_1 + a_1\) , \(s_2 - a_2\) , \ldots, \(s_n + (-1)^{n+1}a_n\) 生成的理想。存在一个 \(\mathbf{A}\) -代数自同构（作用于 \(\mathbf{C}\) ），它把 \(s_k\) 映为 \(s_k + (-1)^{k+1}a_k\) （ \(1 \leqslant k \leqslant n\) ），因而把 \(\mathbf{C}^+\) 映到 \(\mathfrak{c}\) 上；因此我们有 \(\mathbf{C} = \mathbf{A} \oplus \mathfrak{c}\) 。此外，IV第62页定理1的 c) 表明

\[
\mathbf {B} = \oplus \mathbf {C X} ^ {\nu}
\]

，其中 S 是由所有满足下述条件的 \(\nu \in \mathbf{N}^n\) 组成的集合： \(0 \leqslant \nu(i) < i\) （ \(1 \leqslant i \leqslant n\) ）。B 的理想 \(\mathfrak{a}\) 由 \(\mathfrak{c}\) 生成，从而 \(\mathfrak{a} = \mathbf{B}\mathfrak{c} = \bigoplus_{\nu \in \mathbb{S}} \mathfrak{c} \cdot \mathbf{X}^\nu\) 。由于 \(\mathbb{C} = \mathbb{A} \oplus \mathfrak{c}\) ，我们得到

\[
\mathbf {B} = \mathfrak {a} \oplus \bigoplus_ {\nu \in S} \mathbf {A X} ^ {\nu},
\]

由此得出命题5，因为 \(\mathbf{E}_f = \mathbf{B} / \mathfrak{a}\) 。

推论。--- A 到首一多项式 \(f \in A[X]\) 的泛分解代数中的典范同态是单射的。

因为 \(E_{f}\) 的单位元构成 A-模 \(E_{f}\) 的一组基的一部分。

命题6. --- 设 \(f \in \mathbf{A}[\mathbf{X}]\) 是一个次数为 \(n\) 的首一多项式，并设 \(P\) 是 \(\mathbf{X}_1, \ldots, \mathbf{X}_n\) 的一个对称多项式，其系数在 \(A\) 中。则恰好存在一个元素 \(a\) （属于 \(A\) ），它具有以下性质：

(FS) 对任意环同态 \(\rho: \mathbf{A} \to \mathbf{B}\) 以及分解 \({}^{\rho} f(\mathbf{X}) = \prod_{i=1}^{n} (\mathbf{X} - \xi_i)\) （在 \(\mathbf{B}[\mathbf{X}]\) 中），我们有 \(\rho(a) = \mathrm{P}(\xi_1, \ldots, \xi_n)\) 。

记 \(f = \mathbf{X}^n + \sum_{k=1}^{n} a_k \mathbf{X}^{n-k}\) ，则由IV第62页的定理1，存在一个多项式

\(\Pi\) （在 \(n\) 个不定元中，系数在 \(\mathbf{A}\) 中），使得 \(\mathbf{P} = \Pi(s_1, \ldots, s_n)\) 。令 \(a = \Pi(-a_1, a_2, \ldots, (-1)^n a_n)\) 。在假设(FS)下，我们有

\[
s _ {k} (\xi_ {1}, \dots , \xi_ {n}) = (- 1) ^ {k} \rho (a _ {k})
\]

由此

\[
\begin{array}{r l} \rho (a) & = \Pi (- \rho (a _ {1}), \rho (a _ {2}), \dots , (- 1) ^ {n} \rho (a _ {n})) \\ & = \Pi (s _ {1} (\xi_ {1}, \dots , \xi_ {n}), \dots , s _ {n} (\xi_ {1}, \dots , \xi_ {n})) \\ & = P (\xi_ {1}, \dots , \xi_ {n}). \end{array}
\]

这证明了满足(FS)的元素 a 的存在性。a 的唯一性由命题5的推论得出，因为我们有关系 \(a \cdot 1 = P(x_{1,f}, \ldots, x_{n,f})\) （在泛分解代数 \(\mathbf{E}_f\) 中）。

采用命题6的记号，有时我们写 \(a = \mathrm{P}^{\#}(f)\) 。下面是一些例子。

例子。--- 1) 如果 \(\mathrm{P} = s_k\) ，则 \(\mathrm{P}^{\#}(f) = (-1)^k a_k\) 。

2) 设 \(g\) 是 \(\mathbf{A}[\mathbf{X}]\) 中的一个多项式，并令

\[
\mathrm{P} \left(\mathrm{X} _ {1}, \dots , \mathrm{X} _ {n}\right) = g \left(\mathrm{X} _ {1}\right) \dots g \left(\mathrm{X} _ {n}\right).
\]

则 \(\mathrm{P}^{\#}(f)\) 正是结式 \(\operatorname{res}(f, g)\) ，由IV第80页推论1。

\begin{enumerate}
\def\labelenumi{\arabic{enumi})}
\setcounter{enumi}{2}
\item
  令 \(\Delta(X_1, \ldots, X_n) = \prod_{i < j} (X_i - X_j)^2\) ，则 \(\Delta^{\#}(f)\) 正是首一多项式 \(f\) 的判别式（IV，第82页，公式（46））。
\item
  令 \(\mathrm{P}(\mathrm{X}_{1},\ldots,\mathrm{X}_{n})=\mathrm{X}_{1}^{k}+\cdots+\mathrm{X}_{n}^{k}\) ；此外，定义代数 \(\mathrm{E}=\mathrm{A}[\mathrm{X}]/(f)\) ，并用 x 表示 X 在 E 中的像。我们回顾一下，A-模 E 是自由的，其基为 \((1,x,\ldots,x^{n-1})\) （IV，第11页，推论）。让我们证明
\end{enumerate}

\[
\operatorname{Tr} _ {\mathrm{E/A}} (x ^ {k}) = \mathrm{P} ^ {\#} (f).\tag{25}
\]

令 \(\pi_k = \mathrm{Tr}_{\mathrm{E / A}}(x^k)\) （对每个整数 \(k \geqslant 1\) ）。考虑到牛顿关系式（IV，第70页），我们只需建立以下关系式

\begin{enumerate}
\def\labelenumi{(\arabic{enumi})}
\setcounter{enumi}{25}
\tightlist
\item
\end{enumerate}

\[
\begin{array}{l l} \pi_ {k} + a _ {1} \pi_ {k - 1} + \dots + a _ {k - 1} \pi_ {1} + k a _ {k} = 0 & \text { for } \quad 1 \leqslant k \leqslant n \\ \pi_ {k} + a _ {1} \pi_ {k - 1} + \dots + a _ {n - 1} \pi_ {k - n + 1} + a _ {n} \pi_ {k - n} = 0 & \text { for } \quad k > n \end{array}\tag{27}
\]

（我们也把它们称为“牛顿关系式”）。关系式(27)是显然的，因为其左边是以下元素的迹

\[
x ^ {k - n} \left(x ^ {n} + a _ {1} x ^ {n - 1} + \dots + a _ {n - 1} x + a _ {n}\right) = 0.
\]

假设 \(1 \leqslant k \leqslant n\) ，并令

\[
y = x ^ {k} + a _ {1} x ^ {k - 1} + \dots + a _ {k - 1} x + a _ {k} \cdot 1;
\]

，令 \(\mathbf{M} = (m_{ij})\) 为线性映射 \(u\mapsto yu\) （在 E 中）关于基 \((x^{i})_{0\leqslant i\leqslant n-1}\) 的矩阵。我们容易得到关系式

\[
\begin{array}{l l} m _ {i i} = a _ {k} & \text { for } \quad 0 \leqslant i <   n - k \\ m _ {i i} = 0 & \text { for } \quad n - k \leqslant i <   n, \end{array}
\]

由此

\[
\operatorname{Tr} _ {\mathrm{E/A}} (y) = \sum_ {i = 0} ^ {n - 1} m _ {i i} = (n - k) a _ {k}.
\]

此外，我们有

\[
\operatorname{Tr} _ {\mathrm{E/A}} (y) = \pi_ {k} + a _ {1} \pi_ {k - 1} + \dots + a _ {k - 1} \pi_ {1} + n a _ {k},
\]

由此得出公式(26)。

\subsection{结式}

在本节中，我们假设给定了两个正整数 \(p, q\) 和两个多项式 \(f, g\) （在 \(\mathbf{A}[\mathbf{X}]\) 中），形如

\[
\begin{array}{r l} & {f = t _ {p} \mathbf {X} ^ {p} + t _ {p - 1} \mathbf {X} ^ {p - 1} + \dots + t _ {0}} \\ & {g = u _ {q} \mathbf {X} ^ {q} + u _ {q - 1} \mathbf {X} ^ {q - 1} + \dots + u _ {0}} \end{array}
\]

使得 \(\deg f \leqslant p\) ， \(\deg g \leqslant q\) 。对每个整数 \(n \geqslant 0\) ，我们用 \(S_n\) 表示 A{[}X{]} 的子 A-模，它由所有次数 \(< n\) 的多项式组成；它以族 \((X^i)_{0 \leqslant i < n}\) 为基，因此秩为 \(n\) 。

我们给 \(\mathbf{S}_q\times \mathbf{S}_p\) 配备基

\[
\mathrm{B} _ {1} = \left(\left(\mathrm{X} ^ {q - 1}, 0\right), \dots , (\mathrm{X}, 0), (1, 0), \left(0, \mathrm{X} ^ {p - 1}\right), \dots , (0, \mathrm{X}), (0, 1)\right)
\]

以及给 \(\mathbf{S}_{p + q}\) 配备基

\[
\mathbf {B} _ {2} = \left(\mathbf {X} ^ {p + q - 1}, \dots , \mathbf {X}, 1\right).
\]

我们定义线性映射 \(\varphi : S_q \times S_p \to S_{p + q}\) 如下：

\[
\varphi (u, v) = u f + v g
\]

并且我们用 \(\mathbf{M}(f,g,p,q)\) 表示 \(\varphi\) 关于基 \(\mathbf{B}_1\) 与 \(\mathbf{B}_2\) 的矩阵。这是一个阶为 \(p + q\) 的方阵，其指标集合为 \(\{0,1,\dots,p + q - 1\}\) 。它的元素 \(a_{ij}\) 由以下规则给出：

\[
\begin{array}{l l} a) & a _ {i j} = t _ {p - i + j} \quad \text {for} \quad 0 \leqslant j \leqslant q - 1, \\ b) & a _ {i j} = u _ {j - i} \quad \text {for} \quad q \leqslant j \leqslant p + q - 1, \end{array}
\]

，其中 \(t_k\) 取为 0 （如果 \(k \notin [0, p]\) ），且 \(u_k = 0\) （如果 \(k \notin [0, q]\) ）。

例如，当 \(p = 2\) 且 \(q = 3\) 时，我们有矩阵

\[
\left( \begin{array}{c c c c c} t _ {2} & 0 & 0 & u _ {3} & 0 \\ t _ {1} & t _ {2} & 0 & u _ {2} & u _ {3} \\ t _ {0} & t _ {1} & t _ {2} & u _ {1} & u _ {2} \\ 0 & t _ {0} & t _ {1} & u _ {0} & u _ {1} \\ 0 & 0 & t _ {0} & 0 & u _ {0} \end{array} \right).
\]

定义1. --- 在上述记号下，矩阵 \(\mathbf{M}(f,g,p,q)\) 的行列式称为对 \((f,g)\) 关于次数 p 和 q 的结式，或简称为 f 和 g 的结式，如果 \(p = \deg f\) 且 \(q = \deg g\) 。

结式记为 \(\operatorname{res}_{p,q}(f,g)\) ，或简记为 \(\operatorname{res}(f,g)\) （当 \(p=\deg f\) , \(q=\deg g\) ）。

例。--- 1) 给定 A 中的 \(\lambda\) , \(\mu\) ，我们有公式

\[
\begin{array}{l} \operatorname{res} _ {p, 0} (f, \lambda) = \lambda^ {p}, \quad \operatorname{res} _ {0, q} (\mu , g) = \mu^ {q} \\ \operatorname{res} _ {p, 1} (f, \lambda) = \lambda^ {p} t _ {p}, \quad \operatorname{res} _ {1, q} (\mu , g) = (- 1) ^ {q} \mu^ {q} u _ {q}, \end{array}
\]

其证明是直接的。

\begin{enumerate}
\def\labelenumi{\arabic{enumi})}
\setcounter{enumi}{1}
\tightlist
\item
  当 \(p = q = 1\) 时，我们有
\end{enumerate}

\[
\operatorname{res} _ {1, 1} \left(t _ {1} \mathbf {X} + t _ {0}, u _ {1} \mathbf {X} + u _ {0}\right) = t _ {1} u _ {0} - t _ {0} u _ {1}.
\]

注。--- 1) 矩阵 \(\mathbf{M}(g, f, q, p)\) 是由 \(\mathbf{M}(f, g, p, q)\) 经 \(pq\) 次列的对换得到的，因此

\[
\operatorname{res} _ {q, p} (g, f) = (- 1) ^ {p q} \operatorname{res} _ {p, q} (f, g).
\]

\begin{enumerate}
\def\labelenumi{\arabic{enumi})}
\setcounter{enumi}{1}
\tightlist
\item
  设 \(\rho: \mathbf{A} \to \mathbf{B}\) 是一个环同态。定义1直接蕴含公式
\end{enumerate}

\[
\operatorname{res} _ {p, q} \left({}^ {\rho} f, {}^ {\rho} g\right) = \rho \left(\operatorname{res} _ {p, q} (f, g)\right).
\]

\begin{enumerate}
\def\labelenumi{\arabic{enumi})}
\setcounter{enumi}{2}
\tightlist
\item
  给定 \(\lambda, \mu\) （在 A 中），我们有
\end{enumerate}

\[
\operatorname{res} _ {p, q} (\lambda f, \mu g) = \lambda^ {q} \mu^ {p} \operatorname{res} _ {p, q} (f, g).\tag{28}
\]

\begin{enumerate}
\def\labelenumi{\arabic{enumi})}
\setcounter{enumi}{3}
\tightlist
\item
  假设 \(p + q \geqslant 1\) 。根据III第532页公式(28)， \(\varphi\) 的像包含常数多项式 \(\operatorname{res}_{p,q}(f, g)\) 。因此存在一对多项式 \((u, v)\) （其中 \(u \in S_q\) , \(v \in S_p\) ），使得
\end{enumerate}

\[
\operatorname{res} _ {p, q} (f, g) = u f + v g,
\]

由此

\[
\operatorname{res} _ {p, q} (f, g) \in \mathrm{A} \cap (f, g).
\]

这一对 \((u, v)\) 是唯一的，当 \(\operatorname{res}_{p,q}(f, g)\) 在 \(\mathbf{A}\) 中可消去时：因为此时 \(\varphi\) 是单射（III，第524页，命题3）。

\begin{enumerate}
\def\labelenumi{\arabic{enumi})}
\setcounter{enumi}{4}
\tightlist
\item
  假设 \(p \geqslant q\) ，并设 \(h \in \mathbf{A}[\mathbf{X}]\) 是一个次数 \(\leqslant p - q\) 的多项式。让我们证明
\end{enumerate}

\[
\operatorname{res} _ {p, q} (f, g) = \operatorname{res} _ {p, q} (f + g h, g).\tag{29}
\]

因为如果我们写 \(\omega(u, v) = (u, uh + v)\) （对 \((u, v) \in S_q \times S_p\) ），则 \(\omega\) 是 A-模 \(S_q \times S_p\) 的一个自同构，并且我们有

\[
\omega^ {- 1} (u, v) = (u, - u h + v).
\]

\(\omega\) 关于基 \(B_{1}\) 的矩阵是下三角的，其对角元素等于 1。另一方面， \(\varphi \circ \omega\) 把 \((u, v)\) 映为 \(u(f + gh) + vg\) 。现在公式(29)意味着表示 \(\varphi\) 与 \(\varphi \circ \omega\) 的矩阵具有相同的行列式，而这由关系式 det \(\omega = 1\) 得出。

假设 \(f\) 是次数为 \(p\) 的首一多项式；令 \(E = A[X]/(f)\) ，并用 \(x\) 表示 \(X\) 在 \(E\) 中的典范像。我们知道（IV，第11页）， \(E\) 是自由 \(A\) -模，以 \((1, x, ..., x^{p-1})\) 为基。于是我们可以定义范数 \(N_{E/A}(u)\) （对任意元素 \(u\) ，属于 \(E\) ）（III，第543页，定义2）。

命题7. --- 假设 \(f\) 是次数为 \(p\) 的首一多项式；使用上述记号，我们有 \(^{1}\)

\[
\operatorname{res} _ {p, q} (f, g) = \mathrm{N} _ {\mathrm{E/A}} (g (x)).\tag{30}
\]

我们定义 A-线性映射 \(\theta\) ：从 \(S_q \times S_p\) 到 \(S_{p+q}\) ，由 \(\theta(u, v) = uf + v\) 。则 \(\theta\) 把基 \(B_1\) （属于 \(S_q \times S_p\) ）映为序列

\[
(f \mathbf {X} ^ {q - 1}, \dots , f \mathbf {X}, f, \mathbf {X} ^ {p - 1}, \dots , \mathbf {X}, 1)
\]

（其元素属于 \(S_{p+q}\) ）；因此矩阵 \(M_{\theta}\) （即 \(\theta\) 关于基 \(B_{1}\) , \(B_{2}\) 的矩阵）是下三角的，其对角元素等于 1，从而 det \(M_{\theta}=1\) 。

由此可知 \(\theta\) 是双射，并且 \(\operatorname{res}_{p,q}(f,g)\) 等于自同态 \(\varphi' = \varphi \circ \theta^{-1}\) （属于 \(S_{p + q}\) ）的行列式。具体地，我们有

\[
\varphi^ {\prime} (u f + v) = u f + v g\tag{31}
\]

（对任意一对 \((u,v)\) ，属于 \(S_{q}\times S_{p}\) ）。现在我们有 \(A[X]=S_{p+q}+(f)\) 且 \(fS_{q}=S_{p+q}\cap(f)\) ，因此 \(S_{p+q}\) 到 A{[}X{]} 中的典范单射通过取商定义了一个同构 \(\gamma\) （从 \(S_{p+q}/fS_{q}\) 到 E 上）。设 \(\psi\) 为 E 中乘以 \(g(x)\) 的映射。公式 (31) 表明 \(\varphi'\) 在 \(fS_{q}\) 上诱导恒等映射，且 \(\gamma^{-1}\psi\gamma\) 在 \(S_{p+q}/fS_{q}\) 上亦然。于是我们有 det \(\varphi'=\det\psi\) ，从而 (30) 成立，因为 det \(\varphi'=\operatorname{res}_{p,q}(f,g)\) 且 det \(\psi=N_{E/A}(g(x))\) （根据定义）。

推论1. --- 设 \(f \in \mathbf{A}[\mathbf{X}]\) 是一个首一多项式；对每个多项式 \(g \in \mathbf{A}[\mathbf{X}]\) ，以下条件等价：

\begin{enumerate}
\def\labelenumi{(\roman{enumi})}
\item
  res \((f,g)\) 在 A 中可逆；
\item
  存在多项式 \(u, v\) （在 \(\mathbf{A}[\mathbf{X}]\) 中），使得 \(uf + vg = 1\) ；
\item
  \(g(x)\) 在代数 \(\mathbf{A}[\mathbf{X}] / (f)\) 中可逆。
\end{enumerate}

(i) 与 (iii) 的等价性由III第545页命题3及命题7得出；(ii) 与 (iii) 的等价性则是平凡的。

推论2. --- 假设 A 是一个域，并设 \(f, g \in A[X]\) ，则当 \(f, g\) 非零时以下条件等价：

\begin{enumerate}
\def\labelenumi{(\roman{enumi})}
\item
  \(\operatorname{res}(f, g) \neq 0\) ;
\item
  多项式 \(f\) 与 \(g\) 在 \(\mathbf{A}[\mathbf{X}]\) 中互素；
\end{enumerate}

(iii) 对 A 的任意扩域 L，多项式 f 和 g 在 L 中没有公共根。

我们可以立即化简到 \(f\) 是首一的情形（IV，第77页，注记3）。(i) 与 (ii) 的等价性不过是推论1依IV第12页定义1的翻译；(ii) 与 (iii) 的等价性不过是IV第13页推论7。

推论3. --- 对每个 \(\lambda \in A\) ，我们有

\[
\operatorname{res} _ {p, 1} (f, \lambda - X) = f (\lambda), \quad \operatorname{res} _ {1, q} (X - \lambda , g) = g (\lambda).\tag{32}
\]

当 \(f(\mathbf{X}) = \mathbf{X} - \lambda\) 时，代数 E 等于 A，并且我们有 \(x = \lambda\) ；第二个公式 (32) 于是由命题7（IV，第77页）得出。根据注记1和3（IV，第76、77页），我们得出结论

\[
\operatorname{res} _ {p, 1} (f, \lambda - \mathbf {X}) = (- 1) ^ {p} \operatorname{res} _ {1, p} (\lambda - \mathbf {X}, f) = (- 1) ^ {p + p} \operatorname{res} _ {1, p} (\mathbf {X} - \lambda , f) = f (\lambda).
\]

现在假设 f 和 g 都是首一的。我们用 F 表示 A-代数 \(\mathrm{A}[\mathrm{X},\mathrm{Y}]/(f(\mathrm{X}),g(\mathrm{Y}))\) ，并用 x（相应地， y）表示 X（相应地， Y）在 F 中的典范像。

命题8. --- 设 \(f\) 与 \(g\) 分别是次数 \(p\) 与 \(q\) 的首一多项式。在上述记号下，A-模 \(F\) 是自由的，以 \((x^i y^j)_{0 \leq i < p, 0 \leq j < q}\) 为基，并且我们有

\[
\operatorname{res} (f, g) = \mathrm{N} _ {\mathrm{F/A}} (x - y).\tag{33}
\]

令 \(E = A[X]/(f)\) ，且 \(E' = A[Y]/(g)\) 。由II第253页推论1，同态 \(\sigma\) （从 \(E \otimes E'\) 到 F 中，由典范同态 \(A[X] \otimes A[Y] \to A[X,Y]\) 导出）是双射；这就证明了关于 F 的基的断言。我们现在把 E 与它在 F 中在 \(\sigma\) 下的像等同起来。于是，把 Y 映为 y 的、从 \(E[Y]/(g(Y))\) 到 F 中的 E-代数同态是一个同构。

由范数的传递性（III，第546页），我们有

\[
\mathrm{N} _ {\mathrm{F/A}} (x - y) = \mathrm{N} _ {\mathrm{E/A}} (\mathrm{N} _ {\mathrm{F/E}} (x - y)).\tag{34}
\]

由命题7（IV，第77页）， \(\mathbf{N}_{\mathrm{F / E}}(x - y)\) 是多项式 \(g(\mathbf{Y})\) 与 \(x - \mathbf{Y}\) （在 E{[}Y{]} 中）的结式，因此等于 \(g(x)\) （IV，第78页，推论3）。由公式（34）和命题7（IV，第77页），我们因此有

\[
\mathrm{N} _ {\mathrm{F/A}} (x - y) = \mathrm{N} _ {\mathrm{E/A}} (g (x)) = \operatorname{res} (f, g).
\]

命题9. --- 设 \(p_1\) 与 \(q_1\) 为正整数， \(f_1, g_1\) 为 A{[}X{]} 中的多项式，使得 \(\deg f_1 \leqslant p_1\) , \(\deg g_1 \leqslant q_1\) ；于是我们有

\begin{enumerate}
\def\labelenumi{(\arabic{enumi})}
\setcounter{enumi}{34}
\tightlist
\item
\end{enumerate}

\[
\operatorname{res} _ {p, q + q _ {1}} (f, g g _ {1}) = \operatorname{res} _ {p, q} (f, g) \cdot \operatorname{res} _ {p, q _ {1}} (f, g _ {1})\tag{36}
\]

\[
\operatorname{res} _ {p + p _ {1}, q} (f f _ {1}, g) = \operatorname{res} _ {p, q} (f, g) \cdot \operatorname{res} _ {p _ {1}, q} (f _ {1}, g).
\]

我们有 \(\operatorname{res}_{p,q}(f,g)=(-1)^{pq}\operatorname{res}_{q,p}(g,f)\) （IV，第76页）；因此只需证明 (35)。同样，注记3（同处）表明，若公式 (35) 对某个多项式 f 成立，则它对所有形如 \(\lambda f\) （其中 \(\lambda\in A\) ）的多项式也成立。最后，当 f 是 p 次首一多项式时，公式 (35) 由命题7（IV，第77页）及公式 \(\mathrm{N}_{\mathrm{E}/\mathrm{A}}(ab)=\mathrm{N}_{\mathrm{E}/\mathrm{A}}(a)\cdot\mathrm{N}_{\mathrm{E}/\mathrm{A}}(b)\) 得出。总之，我们得出结论：当系数 \(t_{p}\) （即 f 中 \(X^{p}\) 的系数）可逆时，(35) 成立。

引理5. --- 设 t 是 A 的一个元素。存在一个交换环 C（以 A 为子环）、C 的一个子环 B（包含 A）、一个在 C 中可逆的元素 \(\tau\) （属于 B），以及一个环同态 \(\rho: B \to A\) ，使得 \(\rho(\tau) = t\) 且 \(\rho\) 在 A 上的限制等于 \(Id_{A}\) 。

只需对 B 取幺半群 N 的代数 \(\mathbf{A}^{(\mathbf{N})}\) ，即多项式代数 \(A[\tau]\) （单个未定元 \(\tau\) 的），对 C 取群 Z 的代数 \(\mathbf{A}^{(\mathbf{Z})}\) ，对 \(\rho\) 取同态 \(P \mapsto P(t)\) （从 \(\mathbf{A}[\tau]\) 到 A 中）。

采用引理 5 的记号，其中我们已取 \(t = t_p\) ，令

\[
\mathbf {F} = \tau \mathbf {X} ^ {p} + t _ {p - 1} \mathbf {X} ^ {p - 1} + \dots + t _ {1} \mathbf {X} + t _ {0}\tag{37}
\]

在 B{[}X{]} 中。 \(\mathbf{X}^p\) 在 F 中的系数在 C 内可逆；若我们把 F、 \(g\) 、 \(g_1\) 视为 C{[}X{]} 的多项式，则根据前述内容，我们有

\[
\operatorname{res} _ {p, q + q _ {1}} (\mathrm{F}, g g _ {1}) = \operatorname{res} _ {p, q} (\mathrm{F}, g) \cdot \operatorname{res} _ {p, q _ {1}} (\mathrm{F}, g _ {1})\tag{38}
\]

由前述内容。若我们把 F、g 及 \(g_{1}\) 视为 B{[}X{]} 中的多项式，则结式保持不变。由于 \( {}^{\rho}F = f\) , \( {}^{\rho}g = g\) , \( {}^{\rho}g_{1} = g_{1}\) ，公式(35)由(38)及注记2（IV，第77页）得出。

推论1. --- (i) 设 \(\lambda, \alpha_{1}, \ldots, \alpha_{p}\) 为 A 的元素，并假设 \(f(\mathbf{X}) = \lambda(\mathbf{X} - \alpha_{1}) \ldots (\mathbf{X} - \alpha_{p})\) 。我们有

\[
\operatorname{res} _ {p, q} (f, g) = \lambda^ {q} g (\alpha_ {1}) \dots g (\alpha_ {p}).\tag{39}
\]

\begin{enumerate}
\def\labelenumi{(\roman{enumi})}
\setcounter{enumi}{1}
\tightlist
\item
  设 \(\mu, \beta_1, \ldots, \beta_q\) 为 \(A\) 的元素，并进一步假设 \(g(X) = \mu(X - \beta_1) \ldots (X - \beta_q)\) 。于是我们有
\end{enumerate}

\[
\operatorname{res}_{p,q}(f,g) = \lambda^{q}\mu^{p}\prod_{\substack{1\leqslant i\leqslant p\\ 1\leqslant j\leqslant q}}\left(\alpha_{i} - \beta_{j}\right).\tag{40}
\]

断言(i)直接由公式(28)、(32)和(36)得出。现在断言(ii)由(i)得出。

推论2. --- 对每个整数 \(r \geqslant 0\) ，我们有

\[
\operatorname{res} _ {p, q + r} (f, g) = t _ {p} ^ {r} \cdot \operatorname{res} _ {p, q} (f, g).\tag{41}
\]

考虑到例1（IV，第76页），我们只需在 (35) 中取 \(q_{1} = r\) , \(g_{1} = 1\) 。

假设 \(f\) 是首一的，并令 \(\rho: \mathbf{A} \to \mathbf{B}\) 为一个环同态，且 \(\xi_{1}, \ldots, \xi_{p}\) 为 \(\mathbf{B}\) 的元素，使得我们有分解

\[
{ } ^ { p } f ( \mathbf { X } ) = ( \mathbf { X } - \xi _ { 1 } ) \dots ( \mathbf { X } - \xi _ { p } ) .
\]

根据注记2（IV，第77页）及上述推论1，我们有

\[
\rho (\operatorname{res} (f, g)) = g (\xi_ {1}) \dots g (\xi_ {p}).
\]

这一注记特别适用于 \(f(\mathrm{IV}, \mathrm{p.~73})\) 的泛分解，且由于此时 \(\rho\) 是单射的，这提供了一种计算 \(\operatorname{res}(f, g)\) 的方法。

例3. --- 让我们证明公式

\[
\begin{array}{r l} \operatorname{res} _ {2, 2} (a \mathbf {X} ^ {2} + b \mathbf {X} + c, a ^ {\prime} \mathbf {X} ^ {2} + b ^ {\prime} \mathbf {X} + c ^ {\prime}) & = \\ & = (a c ^ {\prime} - c a ^ {\prime}) ^ {2} + (b c ^ {\prime} - c b ^ {\prime}) (b a ^ {\prime} - a b ^ {\prime}). \end{array}\tag{42}
\]

如同命题9（IV，第79页）的论证，我们看到只需证明当a可逆时的公式即可。此时存在如下形式的分解

\[
a \mathbf {X} ^ {2} + b \mathbf {X} + c = a (\mathbf {X} - x) (\mathbf {X} - y)\tag{43}
\]

在 B{[}X{]} 中，其中 B 是一个包含 A 作为子环的适当环。根据上述推论1，所求的结式等于

\[
\mathbf {R} = a ^ {2} \left(a ^ {\prime} x ^ {2} + b ^ {\prime} x + c ^ {\prime}\right) \left(a ^ {\prime} y ^ {2} + b ^ {\prime} y + c ^ {\prime}\right).\tag{44}
\]

现在我们得到

\[
a x + a y = - b, \quad a x y = c
\]

由(43)，因此 \((ax)^2 + (ay)^2 = b^2 - 2ac\) 。

由(44)我们有

\[
\begin{array}{r l} \mathrm{R} & = a ^ {\prime 2} (a x y) ^ {2} + a b ^ {\prime 2} (a x y) + a ^ {2} c ^ {\prime 2} + a ^ {\prime} b ^ {\prime} (a x y) (a x + a y) + \\ & \qquad \qquad \qquad \qquad \qquad \qquad \qquad \qquad \qquad \qquad \qquad \qquad + a ^ {\prime} c ^ {\prime} ((a x) ^ {2} + (a y) ^ {2}) + a b ^ {\prime} c ^ {\prime} (a x + a y) \\ & = a ^ {\prime 2} c ^ {2} + a b ^ {\prime 2} c + a ^ {2} c ^ {\prime 2} - a ^ {\prime} b ^ {\prime} c b + a ^ {\prime} c ^ {\prime} (b ^ {2} - 2 a c) - a b ^ {\prime} c ^ {\prime} b \\ & = (a c ^ {\prime} - c a ^ {\prime}) ^ {2} + (a b ^ {\prime} - a ^ {\prime} b) (b ^ {\prime} c - c ^ {\prime} b), \end{array}
\]

由此所需结果成立。

\subsection{判别式}

定义2. --- 设 f 为 A{[}X{]} 中次数为 m 的首一多项式，并记 E 为 A-代数 A{[}X{]}/(f) ，用 x 表示 X 在 E 中的典范像。我们把 f 的判别式（记为 dis(f)）定义为判别式 \(D_{E/A}(1, x, \ldots, x^{m-1})\) （即 A-代数 E 的基 \((1, x, \ldots, x^{m-1})\) 的判别式）。

对每个正整数 k，我们记 \(p_{k} = \mathrm{Tr}_{\mathrm{E}/\mathrm{A}}(x^{k})\) 。根据III第549页定义2，这等价于公式

\[
\operatorname{dis} (f) = \det \left(p _ {i + j}\right) _ {0 \leqslant i, j <   m}.\tag{45}
\]

例子。--- 1) 如果 \(f\) 是次数为 0 或 1 的首一多项式，则由（45）我们有 \(\operatorname{dis}(f) = 1\) 。

\begin{enumerate}
\def\labelenumi{\arabic{enumi})}
\setcounter{enumi}{1}
\tightlist
\item
  设 \(f(\mathbf{X}) = \mathbf{X}^2 + \alpha \mathbf{X} + \beta\) 为次数为 2 的首一多项式。牛顿关系可以写成如下形式（IV，第75页）
\end{enumerate}

\[
\begin{array}{c} p _ {0} = 2 \\ p _ {1} + \alpha = 0 \\ p _ {2} + \alpha p _ {1} + 2 \beta = 0, \end{array}
\]

从而 \(p_1 = -\alpha\) , \(p_2 = \alpha^2 - 2\beta\) 。由此可得我们有

\[
\operatorname{dis} (f) = \det \left( \begin{array}{c c} 2 & - \alpha \\ - \alpha & \alpha^ {2} - 2 \beta \end{array} \right) = \alpha^ {2} - 4 \beta .
\]

设 B 为一个交换环， \(\rho\) 为从 A 到 B 的同态，且 \(\xi_{1}, \ldots, \xi_{m}\) 为 B 的元素，使得

\[
{ } ^ { \rho } f ( \mathbf { X } ) = ( \mathbf { X } - \xi _ { 1 } ) \dots ( \mathbf { X } - \xi _ { m } ) .
\]

我们用 \(M\) 表示矩阵 \((\rho(p_{i+j}))_{0 \leqslant i,j < m}\) ，用 \(D\) 表示范德蒙德矩阵 \((\xi_{i+1}^{j})_{0 \leqslant i,j < m}\) 。由IV第75页例4，我们有

\[
\rho (p _ {k}) = \xi_ {1} ^ {k} + \dots + \xi_ {m} ^ {k},
\]

从而 \(M = {}^t D \cdot D\) ；由III第532页，我们有 \(\det D = \prod_{i > j} (\xi_i - \xi_j)\) 且 det \(M = (\det D)^2\) ，即

\[
\rho (\operatorname{dis} (f)) = \prod_ {i <   j} \left(\xi_ {i} - \xi_ {j}\right) ^ {2}.\tag{46}
\]

此外（用 D 表示导子 \(\frac{d}{d\mathbf{X}}\) ）我们有

\[
D (^ {\rho} f) (\xi_ {i}) = (\xi_ {i} - \xi_ {1}) \dots (\xi_ {i} - \xi_ {i - 1}) (\xi_ {i} - \xi_ {i + 1}) \dots (\xi_ {i} - \xi_ {m})
\]

对 \(1 \leqslant i \leqslant m\) ，从而

\[
\rho (\operatorname{dis} (f)) = (- 1) ^ {m (m - 1) / 2} \prod_ {i \neq j} (\xi_ {i} - \xi_ {j}) = (- 1) ^ {m (m - 1) / 2} \prod_ {i = 1} ^ {m} D (^ {\rho} f) (\xi_ {i}).
\]

根据IV第80页推论1（应用于 \(f\) 的泛分解），我们最终得到

\[
\operatorname{res} (f, \mathrm{D} f) = \operatorname{res} (\mathrm{D} f, f) = (- 1) ^ {m (m - 1) / 2} \operatorname{dis} (f).\tag{47}
\]

命题10. --- 设 \(m \geqslant 1\) 。存在唯一的多项式 \(\Delta \in \mathbf{Z}[\mathbf{A}_1, \ldots, \mathbf{A}_m]\) ，它具有如下性质：对任意交换环 A 及首一多项式 \(f = \mathbf{X}^m + \sum_{i=1}^{m} a_i \mathbf{X}^{m-i}\) （在 A{[}X{]} 中），我们有

\[
\operatorname{dis} (f) = \Delta \left(a _ {1}, \dots , a _ {m}\right).\tag{48}
\]

此外， \(\Delta\) 的次数 \(\leqslant 2m - 2\) ，并且若赋予权重 \(i\) 给 \(\mathbf{A}_i\) ，则 \(\Delta\) 是等权的，权为 \(m(m - 1)\) 。

\begin{enumerate}
\def\labelenumi{\alph{enumi})}
\item
  \(\Delta\) 的唯一性：若 \(\Delta\) 满足（48），则特别地有 \(\Delta = \text{dis}(F)\) ，其中 \(F\) 是多项式 \(X^{m} + \sum_{i=1}^{m} A_{i} X^{m-i}\) ，其系数在 \(\mathbf{Z}[A_{1}, \ldots, A_{m}]\) 中。
\item
  \(\Delta\) 的存在性：设 \(s_1, \ldots, s_m\) 为未定元 \(\mathbf{X}_1, \ldots, \mathbf{X}_m\) 的初等对称多项式。存在一个多项式 \(\Delta \in \mathbf{Z}[\mathbf{A}_1, \ldots, \mathbf{A}_m]\) ，它是权为 \(m(m-1)\) 的等权多项式，使得
\end{enumerate}

\[
\Delta \left(- s _ {1}, s _ {2}, \dots , (- 1) ^ {m} s _ {m}\right) = \prod_ {i <   j} \left(\mathrm{X} _ {i} - \mathrm{X} _ {j}\right) ^ {2};\tag{49}
\]

因为右边是一个对称多项式 P，它是 \(m(m-1)\) 次齐次的，位于 \(Z[X_{1},\ldots,X_{m}]\) 中（IV，第62页，定理1）。现在，公式（46）表明在IV第74页的记号下有 \(\text{dis}(f)=P(f)\) ，由此直接得出(48)。

\begin{enumerate}
\def\labelenumi{\alph{enumi})}
\setcounter{enumi}{2}
\tightlist
\item
  \(\Delta\) 的次数：关系式(47)与结式的定义（IV，第76页）蕴含公式
\end{enumerate}

\[
(- 1) ^ {m (m - 1) / 2} \Delta = \det \left(a _ {i j}\right) _ {0 \leqslant i, j \leqslant 2 m - 2}\tag{50}
\]

其中 \(a_{ij}\) 取值如下

\[
\begin{array}{l l} a _ {0 0} = 1, \quad a _ {0, m - 1} = m, & a _ {0 j} = 0 \quad \text {if} \quad j \neq 0, \quad j \neq m - 1 \\ a _ {i j} = \mathbf {A} _ {i - j} & \text {for} \quad 1 \leqslant i \leqslant 2 m - 2, \quad 0 \leqslant j \leqslant m - 2 \\ a _ {i j} = (j - i + 1) \mathbf {A} _ {m + i - j - 1} & \text {for} \quad 1 \leqslant i \leqslant 2 m - 2, \quad m - 1 \leqslant j \leqslant 2 m - 2. \end{array}
\]

在这些公式中，约定 \(A_{0}=1\) ，且 \(A_{i}=0\) （对 i\textless0 或 i\textgreater m）。现在(50)立即表明 \(\Delta\) 的次数 \(\leqslant2m-2\) ，这正是我们要证明的。

命题10使我们能够把判别式的定义推广到非首一多项式。设 \(m \geqslant 1\) 为一个整数，则存在唯一的 \(2m - 2\) 次齐次多项式，记作 \(\tilde{\Delta}\) ，它属于 \(\mathbf{Z}[\mathbf{A}_0, \mathbf{A}_1, \ldots, \mathbf{A}_m]\) ，使得

\[
\Delta \left(\mathrm{A} _ {1}, \dots , \mathrm{A} _ {m}\right) = \tilde {\Delta} \left(1, \mathrm{A} _ {1}, \dots , \mathrm{A} _ {m}\right);\tag{51}
\]

这是因为，由于 \(\Delta\) 的次数 \(\leqslant 2m - 2\) ，有理分式

\[
\mathbf {A} _ {0} ^ {2 m - 2} \Delta (\mathbf {A} _ {1} / \mathbf {A} _ {0}, \dots , \mathbf {A} _ {m} / \mathbf {A} _ {0})
\]

属于子环 \(\mathbf{Z}[\mathbf{A}_0, \mathbf{A}_1, \ldots, \mathbf{A}_m]\) （即 \(\mathbf{Q}(\mathbf{A}_0, \mathbf{A}_1, \ldots, \mathbf{A}_m)\) 的子环）。若 \(\mathbf{A}_i\) 的权为 \(i\) （ \(0 \leqslant i \leqslant m\) ），则 \(\tilde{\Delta}\) 是等权的，权为 \(m(m - 1)\) 。若 \(f\) 是次数 \(\leqslant m\) 的多项式，设

\[
f = a _ {0} \mathrm{X} ^ {m} + a _ {1} \mathrm{X} ^ {m - 1} + \dots + a _ {m - 1} \mathrm{X} + a _ {m},
\]

我们令

\[
\operatorname{dis} _ {m} (f) = \tilde {\Delta} \left(a _ {0}, a _ {1}, \dots , a _ {m}\right).\tag{52}
\]

当 \(m = \deg f\) 时，我们将简单地用 \(\operatorname{dis}(f)\) 代替 \(\operatorname{dis}_m(f)\) ；若 \(f\) 是首一的，则由(48)、(51)和(52)， \(\operatorname{dis}(f)\) 与定义2中定义的判别式一致。

命题11. --- 设 \(f\) 为 A{[}X{]} 中次数 \(\leqslant m\) 的多项式。

\begin{enumerate}
\def\labelenumi{(\roman{enumi})}
\item
  若 \(\rho: \mathbf{A} \to \mathbf{B}\) 是一个环同态，则我们有 \(\mathrm{dis}_m({}^{\rho}f) = \rho(\mathrm{dis}_m(f))\) 。
\item
  设 \(\lambda, \alpha_{1}, \ldots, \alpha_{m}\) 为 A 的元素。如果 \(f = \lambda(X - \alpha_{1}) \ldots (X - \alpha_{m})\) ，则我们有
\end{enumerate}

\[
\operatorname{dis} _ {m} (f) = \lambda^ {2 m - 2} \prod_ {i <   j} (\alpha_ {i} - \alpha_ {j}) ^ {2}.\tag{53}
\]

\begin{enumerate}
\def\labelenumi{(\roman{enumi})}
\setcounter{enumi}{2}
\tightlist
\item
  设 \(a_0\) 为 \(\mathbf{X}^m\) 在 \(f\) 中的系数，则我们有
\end{enumerate}

\[
\operatorname{res} _ {m, m - 1} (f, \mathbf {D} f) = \operatorname{res} _ {m - 1, m} (\mathbf {D} f, f) = (- 1) ^ {m (m - 1) / 2} a _ {0} \operatorname{dis} _ {m} (f).\tag{54}
\]

断言(i)是显然的。

由于 \(\tilde{\Delta}\) 是 \(2m - 2\) 次齐次的，我们有

\[
\operatorname{dis} _ {m} (\lambda f) = \lambda^ {2 m - 2} \operatorname{dis} _ {m} (f)\tag{55}
\]

对每个多项式 \(f \in \mathbf{A}[\mathbf{X}]\) （次数 \(\leqslant m\) ）成立。断言(ii)现在由公式(46)和(55)得出。

当 f 是 m 次的首一多项式时，我们有 \(a_{0} = 1\) ，且断言(iii)归结为公式(47)。考虑到(55)以及关系式

\[
\operatorname{res} _ {m, n} (\lambda f, \mu g) = \lambda^ {n} \mu^ {m} \operatorname{res} _ {m, n} (f, g),\tag{56}
\]

（IV，第77页），我们可以从那里过渡到 \(a_0\) 在 A 中可逆的情形。一般情形则由命题11(i)与引理5（IV，第79页）得出。

推论1. --- 设 \(g \in A[X]\) ，并设 \(n\) 为使 \(\deg g \leqslant n\) 成立的正整数。我们有

\[
\operatorname{dis} _ {m + n} (f g) = \operatorname{dis} _ {m} (f) \cdot \operatorname{dis} _ {n} (g) \cdot \operatorname{res} _ {m, n} (f, g) ^ {2}.\tag{57}
\]

通过与前面相同的推理，我们归结为 f 与 g 分别为 m 次和 n 次首一多项式的情形。现在令 \(B = E_{f} \otimes E_{g}\) ，其中 \(E_{f}\) （相应地， \(E_{g}\) ）是 \(f(\text{resp. } g)\) （IV，第73页）的泛分解代数。于是 A 是 B 的一个子环，并且在 B{[}X{]} 中我们有分解

\[
f = \prod_ {i = 1} ^ {m} \left(\mathrm{X} - \alpha_ {i}\right), \quad g = \prod_ {j = 1} ^ {n} \left(\mathrm{X} - \beta_ {j}\right).
\]

因此我们有

\[
f g = \prod_ {k = 1} ^ {m + n} \left(\mathbf {X} - \gamma_ {k}\right),
\]

其中 \(\gamma_{i} = \alpha_{i}\) （ \(1 \leqslant i \leqslant m\) ），且 \(\gamma_{m + j} = \beta_{j}\) （ \(1 \leqslant j \leqslant n\) ）。我们有明显的恒等式

\[
\prod_ {k <   k ^ {\prime}} \left(\gamma_ {k} - \gamma_ {k ^ {\prime}}\right) = \prod_ {i <   i ^ {\prime}} \left(\alpha_ {i} - \alpha_ {i ^ {\prime}}\right) \cdot \prod_ {j <   j ^ {\prime}} \left(\beta_ {j} - \beta_ {j ^ {\prime}}\right) \cdot \prod_ {i, j} \left(\alpha_ {i} - \beta_ {j}\right).
\]

将此关系式平方，并利用(40)和(46)，我们得到(57)。

推论2. --- 若 \(a_{0}\) 是 \(\mathbf{X}^{m}\) 在 \(f\) 中的系数，则我们有

\[
\operatorname{dis} _ {m + 1} (f) = a _ {0} ^ {2} \operatorname{dis} _ {m} (f).\tag{58}
\]

在推论1中取 \(n = 1\) , \(g = 1\) ，并利用公式 \(\operatorname{res}_{m,1}(f,1) = a_0\) （IV，第76页，例1），即得结论。

推论3. --- 设 A 为一个域， f 为 A{[}X{]} 中的非常数多项式。为使 f 与 Df 互素，必要且充分的是 \(\text{dis}(f) \neq 0\) 。

这由命题11(iii)及IV第78页推论2得出。

注记. --- 对上述推论2应用两次可知，我们有 \(\mathrm{dis}_{m}(f)=0\) （对每个次数 \(\leqslant m-2\) 的多项式 f）。

例子. --- 3) 设 \(m = 2\) 。由例2（IV，第81页）我们有 \(\Delta(\mathbf{A}_1, \mathbf{A}_2) = \mathbf{A}_1^2 - 4\mathbf{A}_2\) ，从而 \(\tilde{\Delta}(\mathbf{A}_0, \mathbf{A}_1, \mathbf{A}_2) = \mathbf{A}_1^2 - 4\mathbf{A}_0\mathbf{A}_2\) 。换言之，我们有

\[
\operatorname{dis} _ {2} \left(a _ {0} \mathbf {X} ^ {2} + a _ {1} \mathbf {X} + a _ {2}\right) = a _ {1} ^ {2} - 4 a _ {0} a _ {2}.
\]

\begin{enumerate}
\def\labelenumi{\arabic{enumi})}
\setcounter{enumi}{3}
\tightlist
\item
  考虑多项式
\end{enumerate}

\[
\mathbf {F} = \mathbf {A} _ {0} \mathbf {X} ^ {3} + 3 \mathbf {A} _ {1} \mathbf {X} ^ {2} + 3 \mathbf {A} _ {2} \mathbf {X} + \mathbf {A} _ {3}
\]

其系数在 \(\mathbf{Q}[\mathbf{A}_0,\mathbf{A}_1,\mathbf{A}_2,\mathbf{A}_3]\) 中。我们有

\[
\mathbf {D F} = 3 \left(\mathbf {A} _ {0} \mathbf {X} ^ {2} + 2 \mathbf {A} _ {1} \mathbf {X} + \mathbf {A} _ {2}\right),
\]

\[
\mathrm{F} - 1 / 3 \mathrm{X}. \mathrm{DF} = \mathrm{A} _ {1} \mathrm{X} ^ {2} + 2 \mathrm{A} _ {2} \mathrm{X} + \mathrm{A} _ {3}.
\]

由公式(54)（IV，第84页）和(29)（IV，第77页），我们有

\[
\mathrm{A} _ {0} \cdot \operatorname{dis} _ {3} (\mathrm{F}) = - \operatorname{res} _ {2, 3} (\mathrm{DF}, \mathrm{F}) = - \operatorname{res} _ {2, 3} (\mathrm{DF}, \mathrm{F} - 1 / 3 \mathrm{X}. \mathrm{DF}).
\]

应用IV第80页推论2，我们最终得到

\[
\operatorname{dis} _ {3} (\mathrm{F}) = - 2 7 \operatorname{res} _ {2, 2} \left(\mathrm{A} _ {0} \mathrm{X} ^ {2} + 2 \mathrm{A} _ {1} \mathrm{X} + \mathrm{A} _ {2}, \mathrm{A} _ {1} \mathrm{X} ^ {2} + 2 \mathrm{A} _ {2} \mathrm{X} + \mathrm{A} _ {3}\right).
\]

因此，由IV第80页例3，我们有

\[
\tilde {\Delta} \left(\mathrm{A} _ {0}, 3 \mathrm{A} _ {1}, 3 \mathrm{A} _ {2}, \mathrm{A} _ {3}\right) = - 2 7 \left(\mathrm{A} _ {0} \mathrm{A} _ {3} - \mathrm{A} _ {1} \mathrm{A} _ {2}\right) ^ {2} - 1 0 8 \left(\mathrm{A} _ {1} \mathrm{A} _ {3} - \mathrm{A} _ {2} ^ {2}\right) \left(\mathrm{A} _ {1} ^ {2} - \mathrm{A} _ {0} \mathrm{A} _ {2}\right).
\]

经过一些计算我们发现，如果 \(f = a_0\mathbf{X}^3 + a_1\mathbf{X}^2 + a_2\mathbf{X} + a_3\) ，则我们有

\[
\operatorname{dis} _ {3} (f) = a _ {1} ^ {2} a _ {2} ^ {2} + 1 8 a _ {0} a _ {1} a _ {2} a _ {3} - 4 a _ {1} ^ {3} a _ {3} - 4 a _ {0} a _ {2} ^ {3} - 2 7 a _ {0} ^ {2} a _ {3} ^ {2}.
\]

特别是，我们有

\[
\operatorname{dis} \left(\mathbf {X} ^ {3} + p \mathbf {X} + q\right) = - \left(4 p ^ {3} + 2 7 q ^ {2}\right).
\]

\BourbakiExercises

\BourbakiExerciseGroup

\begin{enumerate}
\def\labelenumi{\arabic{enumi})}
\item
  当 M 是实数加法群 R 的子群时，将一元多项式中的欧几里得除法推广到 M 的群代数上。
\item
  设 \(f\) 为一个多项式 \(\neq 0\) ，它属于 \(\mathbf{A}[\mathbf{X}]\) ，次数为 \(n\) ，首项系数为 \(\alpha_0\) 。设 \(\mathbf{M}\) 为 \(\mathbf{A}[\mathbf{X}]\) （视为 A-模）的子模，由这样的多项式组成：其 \(m\) 次项（ \(m\) 为 \(\mathbf{N}\) 中任意元）的系数可被 \(\alpha_0^\mu\) 整除，其中 \(\mu = \max(m - n + 1, 0)\) ；证明对任意多项式 \(g\) （在 \(\mathbf{M}\) 中），存在两个多项式 \(u, v\) （在 \(\mathbf{A}[\mathbf{X}]\) 中），使得 \(\deg v < n\) 且 \(g = uf + v\) 。
\item
  \begin{enumerate}
  \def\labelenumii{\alph{enumii})}
  \tightlist
  \item
    在 A 上的一元多项式代数 A{[}X{]} 中，映射 \((u,v)\mapsto u(v)\) 是一个内部合成法则；证明该法则满足结合律，并且关于 A{[}X{]} 中的加法与乘法是左分配的。
  \end{enumerate}
\end{enumerate}

\begin{enumerate}
\def\labelenumi{\alph{enumi})}
\setcounter{enumi}{1}
\item
  若 A 是整环，证明关系 \(u(v)=0\) 蕴含 u=0 或 v 为常数；并且当 \(u\neq0\) 且 \(\deg v>0\) 时， \(u(v)\) 的次数等于 u 与 v 的次数之积。
\item
  从现在起假设 A 是一个域。设 u 和 v 是次数 \textgreater0 的两个多项式（在 A{[}X{]} 中）， f 是次数 \textgreater0 的多项式；证明：若 q 是 u 除以 v 的欧几里得除法的商， r 是其余式，则 \(q(f)\) 与 \(r(f)\) 是 \(u(f)\) 除以 \(v(f)\) 的欧几里得除法的商与余式。
\item
  对每个多项式 \(f\) （属于 \(\mathbf{A}[\mathbf{X}]\) ），用 \(\mathrm{I}(f)\) 表示所有形如 \(u(f)\) 的多项式的集合，其中 \(u\) 取遍 \(\mathbf{A}[\mathbf{X}]\) ；它是 \(\mathbf{A}[\mathbf{X}]\) 的一个子环。为使 \(\mathrm{I}(f) = \mathrm{I}(g)\) ，必要且充分的是 \(g = \lambda f + \mu\) ，其中 \(\lambda \neq 0\) 与 \(\mu\) 均属于 \(\mathbf{A}\) 。
\item
  设 \(f\) 与 \(g\) 为两个次数 \(> 0\) 的多项式（在 \(\mathbf{A}[\mathbf{X}]\) 中）；证明 \(I(f) \cap I(g)\) 或者等于 \(\mathbf{A}\) ，或者具有形式 \(I(h)\) ，其中 \(h\) 是次数 \(> 0\) 的多项式（排除第一种可能性，在 \(I(f) \cap I(g)\) 中考虑一个多项式 \(h\) ，其次数 \(> 0\) 且尽可能低）。
\end{enumerate}

\begin{enumerate}
\def\labelenumi{\arabic{enumi})}
\setcounter{enumi}{3}
\tightlist
\item
  \begin{enumerate}
  \def\labelenumii{\alph{enumii})}
  \tightlist
  \item
    假设对每个非零的 \(n \in \mathbf{Z}\) ，关系 \(n\xi = 0\) 蕴含 \(\xi = 0\) （在 \(\mathbf{A}\) 中）。证明映射 \(f \mapsto f(\mathrm{D}_1, \ldots, \mathrm{D}_n)\) 是多项式代数 \(\mathbf{A}[\mathbf{X}_1, \ldots, \mathbf{X}_n] = \mathbf{E}\) 到 \(\mathbf{A}\) 上的自同态代数（即 \(\mathbf{A}\) -模 \(\mathbf{E}\) 的自同态代数）中的单同态（对 \(n\) 用归纳法）。
  \end{enumerate}
\end{enumerate}

\begin{enumerate}
\def\labelenumi{\alph{enumi})}
\setcounter{enumi}{1}
\tightlist
\item
  设 \(m\) 为整数 \(\geqslant 1\) ，使得 \(m: 1 = 0\) （在 \(\mathbf{A}\) 中）。证明 \(\mathrm{D}_i^m = 0\) 对每个偏导数 \(\mathrm{D}_i\) ( \(1 \leqslant i \leqslant n\) ) 在 \(\mathbf{A}[\mathbf{X}_1, \ldots, \mathbf{X}_n]\) 中成立。
\end{enumerate}

\begin{enumerate}
\def\labelenumi{\arabic{enumi})}
\setcounter{enumi}{4}
\tightlist
\item
  设 \(u \in \mathbf{A}[(\mathbf{X}_i)_{i \in I}]\) ， \(r\) 为整数 \(\geqslant 0\) ，并假设在 \(\mathbf{A}\) 中，对每个 \(n > 0\) ，关系 \(n\xi = 0\) 蕴含 \(\xi = 0\) 。如果我们有关系
\end{enumerate}

\[
\sum_ {i \in I} \left(\mathrm{D} _ {i} u\right) (\mathbf {X}) \mathbf {X} _ {i} = r u (\mathbf {X}),
\]

则 \(u\) 是 \(r\) 次齐次的。

\begin{enumerate}
\def\labelenumi{\arabic{enumi})}
\setcounter{enumi}{5}
\tightlist
\item
  设 K 为一个交换域， I 为一个集合， L 为环 K 上由集合 I 生成的自由结合代数（III，第448页，定义2），带有III第458页所定义的分次。我们按照IV第2页的方式定义 L 中元素的次数以及齐次元的概念。a) 证明 \(\deg(uv) = \deg(u) + \deg(v)\) 且 \(\deg(u + v) \leqslant \sup(\deg(u), \deg(v))\) （对 \(u, v \in L\) ）。
\end{enumerate}

\begin{enumerate}
\def\labelenumi{\alph{enumi})}
\setcounter{enumi}{1}
\item
  给定 \(u, v, q, q' \in L\) 满足 \(\deg(u) \geqslant \deg(v)\) 且 \(\deg(qu - q'v) < \deg(ku)\) ；证明存在 \(q'' \in L\) 使得 \(\deg(u - q''v) < \deg(u)\) 。
\item
  给定 \(u, v, q, q' \in L\) ，满足 \(\deg u \geqslant \deg v > \deg (qu - q'v)\) ；证明存在 \(q'' \in L\) 使得 \(\deg (u - q''v) < \deg v\) 。
\item
  给定 \(u, v, u', v' \in L\) ，齐次且满足 \(uv = u'v'\) 及 \(\deg v \geqslant \deg v' \geqslant 0\) ；存在 \(t \in L\) 使得 \(u' = ut, v = tv'\) 。
\item
  给定 \(u, v \in L\) ，齐次，满足 \(uv = vu\) ；存在 \(w \in L\) 、 \(a, b \in K\) 与 \(m, n \in K\) ，使得 \(u = aw^m\) , \(v = bw^n\) 。
\item
  设 \(u, v \in L\) 使得存在 \(a, b \in L\) （非零）满足 \(au = bv\) ；则存在 \(m, d \in L\) ，使得 \(u\) 与 \(v\) 公有的左倍式（相应地，右因式）是 \(m\) 的左倍式（相应地， \(d\) 的右因式）。
\end{enumerate}

\BourbakiExerciseGroup

\begin{enumerate}
\def\labelenumi{\arabic{enumi})}
\item
  假设 A 是一个整环；设 \(f\) 为环 \(\mathbf{A}[X_1, X_2, \ldots, X_n]\) 中次数 \(\leqslant k_i\) （关于 \(X_i\) ，对 \(1 \leqslant i \leqslant n\) ）的一个多项式。对每个 \(i\) （ \(1 \leqslant i \leqslant n\) ）的值，设 \(H_i\) 为 A 的一个含 \(k_i + 1\) 个元素的集合。证明：若 \(f(x_1, x_2, \ldots, x_n) = 0\) 对每个元素 \((x_i) \in \prod_{i=1}^{n} H_i\) 成立，则 \(f = 0\) 。
\item
  假设 A 是一个无限整环；设 \(\Phi\) 为一组多项式 \(\neq0\) （在环 \(A[X_{1},X_{2},\ldots,X_{n}]\) 中）。证明：若 \(\Phi\) 的基数严格小于 A 的基数，则存在 \(A^{n}\) 的一个与 A 等势的子集 H，使得对每个 \(x=(x_{i})\in H\) 及每个 \(f\in\Phi\) ，有 \(f(x)\neq0\) 。
\item
  设 A 是一个整环，并设 B 为 A 的一个无限子集。证明：若多项式 \(f \in A[X]\) 的次数 \textgreater0 ，则 B 在多项式函数 \(x \mapsto f(x)\) 下的像与 B 等势。
\item
  假设 A 的加法群包含一个无限子群 G，其非零元素在 A 中不是零因子。证明映射 \(f \mapsto \tilde{f}\) ——它把 \(A{[}X_{1}, \ldots, X_{p}{]}\) 映入 \(A^{p}\) 到 A 的映射代数——是单射（注意，单未定元的一个 n 次多项式在 G 中不可能有多于 n 个不同的根）。特别地，当 A 中存在一个元素 \(x_{0}\) ——它不是零因子且在 A 的加法群中是无限阶的——时，此结论成立。
\item
  设 K 是一个具有无穷多个元素的域，并设 Q 为 K 上的四元数代数，类型为 \((-1,0,-1)\) （III，第445页）。证明多项式 \(X^{2}+1\) 在 Q 中有无穷多个根。
\item
  设K是一个含有q个元素的有限域。
\end{enumerate}

\begin{enumerate}
\def\labelenumi{\alph{enumi})}
\tightlist
\item
  在环 \(\mathbf{K}[\mathbf{X}_1, \mathbf{X}_2, \ldots, \mathbf{X}_n]\) 中，令 \(\mathfrak{a}\) 为由 \(n\) 个多项式
\end{enumerate}

\(X_{i}^{q}-X_{i}\quad(1\leqslant i\leqslant n)\) 所生成的理想。证明若 \(f\in\mathfrak{a}\) ，则 \(f(x_{1},x_{2},\ldots,x_{n})=0\) 对所有 \((x_{i})\in\mathbf{K}^{n}\) 成立（注意乘法群 \(K^{*}\) 的阶为 q-1）。

\begin{enumerate}
\def\labelenumi{\alph{enumi})}
\setcounter{enumi}{1}
\item
  若 \(f\) 为 \(\mathbf{K}[\mathbf{X}_1, \mathbf{X}_2, \ldots, \mathbf{X}_n]\) 中任意多项式，证明存在恰好一个多项式 \(\vec{f}\) ，使得 \(\deg_i \vec{f} \leqslant q - 1\) （对 \(1 \leqslant i \leqslant n\) ），且 \(f \equiv \vec{f} \pmod{\mathfrak{a}}\) 。于是有 \(\deg \vec{f} \leqslant \deg f\) ；若 \(f\) 使得 \(f(x_1, x_2, \ldots, x_n) = 0\) （对每个 \((x_i) \in \mathbf{K}^n\) ），则 \(f\) 属于理想 \(\mathfrak{a}\) ，因而该理想是同态 \(f \mapsto \vec{f}\) 下 0 的原像（利用习题1）。
\item
  设 \(f_1, f_2, \ldots, f_m\) 为 \(\mathbf{K}[\mathbf{X}_1, \mathbf{X}_2, \ldots, \mathbf{X}_n]\) 中的非零多项式，使得 \(f_i(0, 0, \ldots, 0) = 0 (1 \leqslant i \leqslant m)\) ，且这些 \(f_i\) 的总次数之和 \(< n\) 。证明存在一个元素 \((x_1, x_2, \ldots, x_n) \in \mathbf{K}^n\) ，它异于 \((0, 0, \ldots, 0)\) ，使得
\end{enumerate}

\[
f _ {i} (x _ {1}, x _ {2}, \dots , x _ {n}) = 0 \quad \text { for } \quad 1 \leqslant i \leqslant m.
\]

（观察到如果并非如此，则多项式 \(\prod_{i=1}^{m} (1 - f_i^{q-1})\) 将模 a 与多项式 \(\prod_{i=1}^{n} (1 - X_i^{q-1})\) 同余，并利用 \(b\) ).

\begin{enumerate}
\def\labelenumi{\arabic{enumi})}
\setcounter{enumi}{6}
\item
  设 K 为含 q 个元素的有限域， B 为乘积环 \(K^{1}\) ，其中 I 是一个无限集。给出 B{[}X{]} 中一个非零多项式 f 的例子，使得 \(f(x)=0\) 对所有 \(x\in B\) 成立。
\item
  设 B 为加法群 \(\mathbf{Z} \oplus (\mathbf{Z}/2\mathbf{Z})\) ，并赋予它由令 \((a, b) \cdot (a', b') = (aa', ab' + ba')\) （对 \((a, b) \in \mathbf{B}\) 与 \((a', b') \in \mathbf{B}\) ）所得到的交换环结构。设 \(\varepsilon = (0, 1) \in \mathbf{B}\) ，以及 \(f(\mathbf{X}) = \varepsilon \mathbf{X}(\mathbf{X} - 1) \in \mathbf{B}[\mathbf{X}]\) 。则 \(f(\alpha) = 0\) 对所有 \(\alpha \in B\) 成立。
\item
  设 \(k\) 为整数 \(> 0\) ，使得 \(k \cdot 1 = 0\) （在 \(\mathbf{A}\) 中）。设 \(h\) 为整数 \(> 0\) ， \(f\) 为多项式 \((\mathbf{X} - \alpha)^{k + h} + (\mathbf{X} - \alpha)^k\) 。则 \(\alpha\) 是阶为 \(k\) 的根（对 \(f\) 而言），且是阶 \(\geqslant k + h - 1\) 的根（对 \(Df\) 而言）。
\end{enumerate}

10) 设 \(f\) 为 \(\mathbf{Z}[\mathbf{X}_1, \ldots, \mathbf{X}_n]\) 的一个非零元素。设 \(\mathrm{N}(q)\) 为 \(f\) 在胞腔 \(|x_1| \leqslant q, \ldots, |x_n| \leqslant q\) 中的零点个数。则 \(\mathrm{N}(q) / q^n\) 当 \(q \to +\infty\) 时趋于 0 。

\begin{enumerate}
\def\labelenumi{\arabic{enumi})}
\setcounter{enumi}{10}
\item
  设 P 为所有素数的集合， \(P'\) 为 P 的一个无限子集，且 \(f \in \mathbf{Z}[(\mathbf{X}_i)]\) 。对每个 \(p \in P'\) ，令 \(f_p\) 为 \((\mathbf{Z}/p\mathbf{Z})[(\mathbf{X}_i)]\) 中由对 f 的系数作用 Z 到 Z/pZ 上的典范同态所得到的元素。假设对每个 \(p \in P'\) 及每个 \(x \in (\mathbf{Z}/p\mathbf{Z})^l\) ，我们有 \(f_p(x) = 0\) ；则 f = 0。
\item
  设 \(\mathbf{M}\) 为一个 A-模，且 \(\alpha_0, \ldots, \alpha_n \in \mathbf{A}\) 使得：对 \(i \neq j\) ，关于 \(\alpha_i - \alpha_j\) 的倍乘同态在 \(\mathbf{M}\) 上是双射。对任意 \(m_0, \ldots, m_n \in M\) ，存在且仅存在一个 \(f \in \mathbf{M}[\mathbf{X}]\) ，使得 \(f(\alpha_i) = m_i\) （对 \(0 \leqslant i \leqslant n\) ）且 \(\deg f \leqslant n\) 。
\item
  设 \(\alpha_{i}\) ( \(1 \leqslant i \leqslant n\) ) 为 \(n\) 个互异元素（属于一个交换域 \(\mathbf{K}\) ）， \(\beta_{i}\) ( \(1 \leqslant i \leqslant n\) ) 与 \(\gamma_{i}\) ( \(1 \leqslant i \leqslant n\) ) 为 \(2n\) 个任意元素（属于 \(\mathbf{K}\) ）。证明：存在且仅存在一个多项式 \(f \in \mathbf{K}[\mathbf{X}]\) ，次数 \(\leqslant 2n - 1\) ，使得 \(f(\alpha_{i}) = \beta_{i}\) 且 \(f'(\alpha_{i}) = \gamma_{i}\) （对 \(1 \leqslant i \leqslant n\) ）（埃尔米特插值公式）（首先考虑 \(2n - 1\) 个（在 \(2n\) 个元素 \(\beta_{i}, \gamma_{i}\) 中）为零的情形）。
\end{enumerate}

\BourbakiExerciseGroup

\begin{enumerate}
\def\labelenumi{\arabic{enumi})}
\item
  设 E 是交换域 K 上的一个结合、交换且保单位的代数；设 \(\mathbf{x} = (x_i)_{i \in I}\) 是 E 中两两可交换的元素族。设 U 是 \(\mathrm{K}(\mathrm{X}_i)_{i \in I}\) 的子环，由所有 \(f \in \mathrm{K}(\mathrm{X}_i)_{i \in I}\) （x 可代入 f）组成。证明：若 E 中非可逆元组成的集合是一个理想，则 U 中亦然。并证明此时 U 中非可逆元组成的理想是极大的。
\item
  \begin{enumerate}
  \def\labelenumii{\alph{enumii})}
  \tightlist
  \item
    设 K 为一个交换域，且 \(u = a_{m}X^{m} + a_{m+1}X^{m+1} + \cdots + a_{n}X^{n}\) 为 K{[}X{]} 中的多项式，使得 \(a_{m} \neq 0\) , \(a_{n} \neq 0\) ( \(0 \leq m \leq n\) )。证明：若 g 是 K(X) 中次数为 \(d \neq 0\) 的有理分式，则 \(u(g)\) 非零，且当 d \textgreater{} 0 时次数为 nd，当 d \textless{} 0 时次数为 md。
  \end{enumerate}
\end{enumerate}

\begin{enumerate}
\def\labelenumi{\alph{enumi})}
\setcounter{enumi}{1}
\tightlist
\item
  由此推断：K(X) 的每个非常数有理分式 g 都可代入 K(X) 中的每个有理分式（注意，若 g 的次数为 0，则存在 \(\alpha \in K\) 使得 \(g - \alpha\) 的次数 \textless{} 0）。
\end{enumerate}

\begin{enumerate}
\def\labelenumi{\arabic{enumi})}
\setcounter{enumi}{2}
\tightlist
\item
  一个有理分式 \(f \in K(X_1, X_2, \ldots, X_n)\) 被称为齐次的，如果它等于两个齐次多项式之商（分母 \(\neq 0\) ). 证明：为使 \(f\) 是齐次的，必要且充分条件是
\end{enumerate}

\[
f \left(\mathrm{ZX} _ {1}, \mathrm{ZX} _ {2}, \dots , \mathrm{ZX} _ {n}\right) = \mathrm{Z} ^ {d} f \left(\mathrm{X} _ {1}, \dots , \mathrm{X} _ {n}\right)
\]

其中 \(d\) 是 \(f\) 的次数。

\begin{enumerate}
\def\labelenumi{\arabic{enumi})}
\setcounter{enumi}{3}
\tightlist
\item
  证明：拉格朗日插值公式（IV，第16页，命题6）可以写成
\end{enumerate}

\[
f (\mathbf {X}) = \omega (\mathbf {X}) \sum_ {i = 1} ^ {n} \frac {\beta_ {i}}{\omega^ {\prime} (\alpha_ {i}) (\mathbf {X} - \alpha_ {i})}
\]

其中 \(\omega(X) = (X - \alpha_1)(X - \alpha_2) \ldots (X - \alpha_n)\) 。由此推出，如果 \(g\) 是 \(K[X]\) 中次数 \(\leqslant n - 2\) 的多项式，则

\[
\sum_ {i = 1} ^ {n} \frac {g (\alpha_ {i})}{\omega^ {\prime} (\alpha_ {i})} = 0
\]

（考虑多项式 \(f(\mathbf{X}) = \mathbf{X}g(\mathbf{X})\) ）。

5) 设 \(\mathbf{K}\) 为特征 0 的交换域。证明：若有理分式 \(u \in \mathbf{K}(\mathbf{X})\) 使得 \(\mathbf{D}u = 0\) ，则 \(u\) 是常数。

\begin{enumerate}
\def\labelenumi{\arabic{enumi})}
\setcounter{enumi}{5}
\item
  将欧拉恒等式推广到特征 0 的域上的齐次有理分式（习题3）（考虑有理分式 \(\frac{1}{Z^m} f(ZX_1, ZX_2, ..., ZX_n)\) 关于 \(Z\) ，并利用习题5）。
\item
  设 \(\mathbf{K}\) 为一个交换域。证明：
\end{enumerate}

\[
[ \mathbf {K} (\mathrm{T}): \mathbf {K} ] = \operatorname{Sup} (\operatorname{Card} (\mathbf {K}), \operatorname{Card} (\mathbf {N})).
\]

(利用 \(1 / (T - a)\) （ \(a \in \mathbf{K}\) ）的线性无关性，来表明

\[
[ \mathrm{K} (\mathrm{T}): \mathrm{K} ] \geqslant \operatorname{Card} (\mathrm{K}).)
\]
